% Modernised reading — fonds Alexandre Grothendieck, shelfmark 68
% (pages 1-55, the whole folder).
%
% This is an INTERPRETATION, not a transcription. It re-reads the pages in
% today's notation and vocabulary, states the hypotheses the manuscript leaves
% implicit, and drops the critical apparatus so that the mathematics reads
% straight through. Everything the brackets and underlines carried moves into a
% footnote. It is held to being mathematically correct as it stands; where the
% manuscript is loose or mistaken, the true statement is what appears here and
% the footnote says what the page has. In any dispute of substance the
% transcription governs: whoever cites Grothendieck must cite
% batch-01.fr.tex (pp. 1-20), batch-02.fr.tex (pp. 21-40) or
% batch-03.fr.tex (pp. 41-55).
%
% Pass: Opus 5.5 (claude-opus-5-5), 2026-09-22 — first pass, unchecked against the pages by a human.
% The folder was read whole, its three batches together; this is one document
% for the shelfmark. It was made on Opus 5.5 and has not been measured against
% the reference reading of folder 115 (made on Opus 5). The references in the
% footnotes are given from memory and were not checked against the sources.
%
% Revised 2026-09-23 (Opus 5.5), after batch 1 (pp. 16, 19, 20) and batch 2
% (p. 21) were re-read on the facsimile: « prétassé » and « tassé » are now
% the pages' own words (batch 1 had « préféré », « bornés »), and the section
% IV footnotes say so; p. 20's « Cela implique que » (was « C'est impliqué
% par ») makes tassé => prétassé the page's own statement, so the departure
% is withdrawn; a footnote on the numeral (4 and 5 superposed) heading p.
% 20; the p. 16 footnote now says the page has [-i, i] and E_1 (sense E_0);
% the p. 52 footnote says the « 82 » is on that copy alone, is not
% transcribed, and is in a hand the facsimile does not attribute.
%
% Spine. Five subjects on loose leaves; no link between them is written.
%  (1) Positional games, in three moutures kept distinct:
%      A  pp. 1-13, pencil, numbered 1-5 (6 is a title only): data
%         (C, R, J, i, G_0), parties, C_i, strategies, G(j), Thm 1, the
%         largest strategy, acceptable positions, Thm 2, J(j'), G_i(j) and
%         « ? Th 3 »;
%      B  pp. 23-33, « Jeux de position »: the fixed-point theorem with a
%         complete proof (well-ordering), strategies against j, admissible
%         positions, the two-player partition, alternating games (G, P, N),
%         the characterisation of (G, P) without infinite parties;
%      C  p. 47, « Jeux à N joueurs »: data and circular games only.
%  (2) Rearrangements on a group: scratch pp. 14-16; numbered run
%      pp. 18-21 (|f,g|, weak majorisation, couples prétassés and tassés,
%      condition K, the class S~).
%  (3) « Exponentielles (en car. 0) », pp. 35-45.
%  (4) Orientation of real vector spaces, pp. 49, 51, 53.
%  (5) Graded determinants, p. 55 (axioms only).
%  Not mathematics, not transcribed: p. 17 (typed administrative table),
%  pp. 50, 52, 54 (typescript « XVIII Les épousailles (2) »). Blank: 22, 34,
%  46, 48.
%
% Notation fixed for the document.
%  - Games: C, R, C_0, J, G_0(j), G(j) as on the pages. The initiative map is
%    alpha (B's letter) in both moutures; A writes i. Cbar(j) is
%    (C \ C_0) \ C(j) throughout (B's convention; A's Cbar(j) = C \ C(j)
%    contains C_0 and A writes Cbar(j) \ Cbar_0(j)). A's « acceptable » and
%    B's « admissible » are kept, each in its mouture. F_j (the operator of
%    the fixed-point equation) and G_alpha(j) for ordinal alpha are ours.
%    B's G, P, N (gagnantes, perdantes, nulles) are kept; combinatorial game
%    theory's N-, P-, D-positions are his G, P, N — letter clash flagged.
%  - Rearrangements: f^down for the decreasing sequence the p. 18 margin
%    calls A(f); e for the neutral element (the pages write 0); prétassé /
%    tassé as the pages name them (pp. 19-21).
%  - Exponentials: Ghat_a for the page's G_a barre; W(U) for K' -> U (x) K';
%    O_X for his underlined O_X.
%  - Orientation: tau_P for the page's underlined a_P; D(W,V) for his
%    Dem_{W,V} / Omega(W,V) (the two half-spaces).
%  - Determinants: Pic^Z(k) for his P(k) with a doubled stem.
%
% Substantive departures, with pages.
%  p. 2    « no infinite partie » does not imply C = union C_i when C is
%          infinite; ordinal rank supplied in a footnote (ours).
%  pp. 3-6 Thm 1 stated in general; A's proof of 2° <= (p. 5) holds only
%          without infinite parties (counterexample ours); p. 6's attempt at
%          the general case is cancelled and incomplete; B pp. 26-28 proves it.
%  pp. 5-8 corollary: « Gamma stable » defined here (ours, from p. 7's
%          struck attempts); valid without infinite parties; no largest
%          strategy in general (counterexample ours).
%  p. 9    G(j) read G_0(j').
%  p. 12   « ? Th 3 »: true if C = union C_i (the page's proof), if C is
%          finite, or if opponents' positions have finitely many successors
%          (ours); false in general, even without infinite parties, so the
%          margin's « pas de partie infinie » is not enough (counterexample
%          ours); in general G(j) = union over ordinals of G_alpha(j), the
%          least fixed point of F_j (ours).
%  p. 25   B's theorem: G(j) is a fixed point, the least one (ours); the
%          equations do not characterise it when there are infinite parties.
%  p. 28   « jouée suivant x_0 » read « suivant Sigma' ».
%  p. 33   condition 1) restricted to Gtilde \ C_0 (as written it forces
%          Gtilde cap C_0 = empty and makes 3) void).
%  p. 18   E_alpha(g') read E_beta(g'); f, g read f', g'; c) taken for all
%          n >= 1 with sums over i < n (the page's n in N*, sum from 0 to n,
%          misses the first term).
%  p. 19   (f*g)(s) = <f * eps_s, gcheck> is (f*g)(s^-1); written
%          <f, eps_s * gcheck>.
%  p. 20   tassé => prétassé, as the page says; the converse fails
%          (counterexample ours).
%  p. 21   a) f = eps_t * f read fcheck = eps_t * f.
%  p. 35   « exp 1 = 0 » read log 1 = 0; the margin's « réciproque ? »
%          answered (ours). Ex. 2 needs J nilpotent or M of finite type
%          (counterexample ours).
%  p. 38   (**) for M not of finite type: Hom(G_a, Aut_K(M)) corresponds to
%          locally nilpotent endomorphisms, not N(End_K(M)); « u_n nuls sauf
%          un nombre fini » fails (counterexample ours).
%  p. 40   c_ij = -1/(i+j+1)! (page 1/(i+j)!); Theta, Theta' act on
%          Hom_K(M, N), not End_K(M).
%  p. 43   E_{Theta+Theta'} = E_Theta o E_Theta' (page +); « loc.
%          nilpotente » read as nilpotent locally on X (for a locally
%          nilpotent derivation in today's sense the statement fails).
%  p. 44   Cor: (i)-(iii) do not imply (iv); exp maps nilpotent derivations
%          onto unipotent automorphisms, a strictly smaller set
%          (counterexample ours; cf. the page's own « est-il nécessaire ? »).
%  p. 45   last Cor: a), b) false as stated; true with Theta(I) in I^2 and
%          u = id on gr_I(A) (the struck margin (iii')); counterexample ours.
%  p. 49   St_V read St_{V'}; Omega(V,W) wedge Omega(W) read ... Omega(V).
%  p. 53   the sign « (-1)^? » fixed in p. 51's normalisation (ours).
%  p. 55   phi_{W1,W1} read phi_{W1,W2}.
%
% Announced and never established: A's §6 (p. 13, title only); « ? Th 3 » in
% the generality the margin claims; the end of p. 21 (the common point of the
% level sets); any use of p. 16's axioms; the NB question of p. 43; I(omega,
% omega') on p. 53; the theorem p. 55's axioms prepare; any winning condition
% for the N-player games of p. 47.
\documentclass[11pt,a4paper]{article}
% Le préambule commun à toutes les transcriptions du fonds Grothendieck.
%
% Il tient en une idée : **l'incertitude se compose, elle ne se résout pas.**
% Une transcription qui lisse un mot douteux a détruit la seule chose pour
% laquelle on la faisait — un lecteur doit pouvoir savoir, à chaque endroit,
% ce qui était sur le feuillet et ce qui a été ajouté. Les six macros qui
% suivent sont donc l'appareil critique, et non de la décoration.
%
% Les mêmes commandes sont reconnues par `scripts/render.mjs`, qui produit la
% vue de lecture du site. Une macro ajoutée ici doit l'être aussi là-bas,
% faute de quoi le rendu échouera bruyamment — ce qui est voulu : un rendu
% silencieusement faux est pire qu'un rendu absent.

\ProvidesPackage{grothendieck}[2026/08/08 Transcriptions du fonds Grothendieck]

\usepackage{amsmath,amssymb,amsthm}
\usepackage{mathrsfs}
\usepackage{stmaryrd}
% Les diagrammes commutatifs sont partout dans ce fonds : c'est la langue
% même de la géométrie algébrique de Grothendieck, et une transcription qui
% les rendrait en prose ne transcrirait rien.
\usepackage{tikz-cd}
% Français par défaut — c'est la langue des feuillets — mais l'anglais est
% chargé aussi : la traduction et le résumé sont des documents anglais, et la
% typographie française (espaces avant les deux-points) y serait une faute.
% Ces documents basculent par \selectlanguage{english} en tête de corps.
\usepackage[english,french]{babel}
\usepackage{fontspec}
\usepackage{geometry}
\geometry{a4paper,margin=25mm,marginparwidth=28mm}
\usepackage{marginnote}
\usepackage{xcolor}
% ulem, not soul. `soul`'s \st and \ul are built on a character-by-character
% scan that XeLaTeX's font machinery breaks: the document fails with an
% undefined control sequence rather than degrading. ulem does the same two jobs
% and is Unicode-engine safe. `normalem` stops it hijacking \emph, which the
% transcriptions use for Grothendieck's own emphasis.
\usepackage[normalem]{ulem}
\usepackage{hyperref}
\hypersetup{colorlinks=true,linkcolor=black,urlcolor=black,pdfborder={0 0 0}}

% --- Métadonnées du lot -----------------------------------------------------
% Reprises telles quelles par le script de rendu, qui les affiche en tête de
% la vue de lecture. Elles fixent ce que le fichier prétend couvrir : sans
% elles, une transcription flotte sans point d'attache dans un fonds de seize
% mille feuillets.

\newcommand{\folder}[1]{\gdef\grothFolder{#1}}
\newcommand{\batch}[1]{\gdef\grothBatch{#1}}
\newcommand{\leaves}[2]{\gdef\grothFirst{#1}\gdef\grothLast{#2}}
\newcommand{\foldertitle}[1]{\gdef\grothFoldertitle{#1}}
\gdef\grothFolder{?}\gdef\grothBatch{?}\gdef\grothFirst{?}\gdef\grothLast{?}\gdef\grothFoldertitle{}

% --- Le résumé --------------------------------------------------------------
%
% Il ouvre la lecture modernisée, sous le titre « Résumé », et il est composé
% en petit. Ce n'est pas un ornement : c'est le seul passage du document qui
% ne soit pas des mathématiques, et un lecteur doit pouvoir le reconnaître
% avant de l'avoir lu — puis le sauter s'il connaît déjà le sujet, ou s'y
% arrêter s'il ne le connaît pas. Le filet qui le suit marque la reprise du
% corps.

\newenvironment{resume}
  {\small\setlength{\parskip}{0.45em}}
  {\par\medskip\hrule\medskip}

% --- Mots-clés ---------------------------------------------------------------
%
% En anglais, en clôture du résumé de la lecture modernisée : le vocabulaire
% moderne sous lequel chercher ce que les feuillets construisent. Le manifeste
% du site (`npm run manifest`) les extrait pour étiqueter la cote entière —
% cette ligne est la source unique des tags, il n'y a pas de fichier à côté.
\newcommand{\keywords}[1]{\par\smallskip\noindent
  {\footnotesize\textsc{Keywords} --- \textit{#1}}\par}

% --- L'appareil critique ----------------------------------------------------

% Le numéro de feuillet, en marge. C'est l'ancre qui permet de régler un
% désaccord sur une formule en regardant *un* feuillet plutôt qu'une chemise
% de six cents. Le site s'en sert aussi pour tourner les pages du fac-similé
% quand on fait défiler la transcription.
\newcommand{\leaf}[1]{%
  \marginnote{\footnotesize\color{gray!70}\textsf{#1}}%
  \label{leaf:#1}\ignorespaces}

% Illisible : on ne devine pas. Un mot inventé qui se lit comme les autres est
% le pire résultat possible.
\newcommand{\ill}{\textcolor{red!60!black}{[\,\ldots\,]}}

% Lecture douteuse : le mot est donné, et signalé comme incertain.
\newcommand{\uncertain}[1]{\uline{#1}}

% Ajout de l'éditeur — une lettre manquante, un mot sous-entendu.
\newcommand{\add}[1]{\textcolor{blue!50!black}{[#1]}}

% Biffé par Grothendieck lui-même. Ce qu'il a rayé fait partie du document :
% une rature dit souvent où la pensée a changé de direction.
\newcommand{\struck}[1]{\sout{#1}}

% Note du transcripteur, sur l'état matériel du feuillet.
\newcommand{\note}[1]{\par\smallskip{\small\color{gray!60!black}\textit{[#1]}}\par\smallskip}

% Note marginale de Grothendieck, distincte du corps du texte.
\newcommand{\marginal}[1]{\par\smallskip\noindent\hspace{1em}%
  {\small\color{gray!60!black}#1}\par\smallskip}

% --- Titre ------------------------------------------------------------------

\AtBeginDocument{%
  \begin{center}
    {\large\bfseries Fonds Alexandre Grothendieck --- cote n\textsuperscript{o}~\grothFolder}\\[2pt]
    {\itshape\grothFoldertitle}\\[4pt]
    {\small feuillets \grothFirst--\grothLast\ (lot \grothBatch)}\\[2pt]
    {\footnotesize\color{gray!60!black}Transcription établie d'après le fac-similé numérisé
     par l'Université de Montpellier.\\
     Les lectures incertaines sont soulignées, les illisibles notées [\,\ldots\,],
     les ajouts entre crochets.}
  \end{center}
  \vspace{1em}}

\endinput


\folder{68}
\pages{1}{55}
\dating{[à partir de 1978-à partir de 1983]}
\watermark{Édition de démonstration\\interprétation personnelle de l'œuvre}
\foldertitle{Jeux de position : notes manuscrites (s.d.) — lecture modernisée du dossier entier}

\begin{document}

\section*{Résumé}

\begin{resume}
Cinquante-cinq pages de notes, au crayon et à l'encre bleue, sans date de
sa main ; l'inventaire les date d'une période qui s'ouvre en 1978 et se ferme
en 1983 ou après, et les range dans un groupe de dossiers intitulé « Autour
de l'enseignement ». La chemise porte le titre
des premières pages, « Jeux de position », mais les feuillets traitent de cinq
sujets que rien sur les pages ne relie.

Le premier, et le plus développé, est celui des jeux. Pensez aux échecs, ou à
un pion qu'on pousse de case en case : des positions, des coups permis d'une
position à l'autre, des joueurs dont c'est tour à tour le tour, et des
positions finales où l'un a gagné. La question est de savoir de quelles
positions un joueur est sûr de gagner, quoi que fassent les autres. L'idée est
de raisonner à rebours depuis la fin. Une position où c'est à moi de jouer est
gagnante si l'un de mes coups mène à une position gagnante ; une position où
c'est aux autres de jouer est gagnante si tous leurs coups y mènent. On part
des positions finales gagnées et l'on remonte, comme l'eau qui gagnerait un
terrain depuis ses points les plus bas : gagnant en un coup, en deux coups, et
ainsi de suite. Tout le délicat est dans les parties qui peuvent ne pas
finir. Grothendieck écrit deux fois la théorie. La première rédaction a une
démonstration qui ne passe que si toute partie s'achève, et se termine sur un
énoncé précédé d'un point d'interrogation : la remontée coup par coup
atteint-elle toutes les positions gagnantes ? La réponse est non en général :
il faut parfois la poursuivre au-delà de tous les entiers. La seconde
rédaction repart de zéro, démontre le théorème central en toute généralité par
un argument de bon ordre, et traite le jeu à deux joueurs qui jouent
alternativement, où chaque position est gagnante, perdante ou nulle. Une
troisième, d'une page, pose le cadre d'un jeu à $N$ joueurs.

Le deuxième sujet est un problème d'arrangement. Deux listes de nombres
positifs qu'on multiplie terme à terme et qu'on additionne donnent le plus
grand total quand on range le plus grand en face du plus grand. Sur un groupe,
où l'on combine deux fonctions par convolution, la question devient : comment
disposer les valeurs de chaque fonction pour que la convolution ait le plus
haut pic, ou mieux, domine toutes les autres dispositions ? La réponse
attendue est de \emph{tasser} chaque fonction autour d'un centre, en tas
symétrique, comme on tasse du sable ; les pages cherchent sous quelles
conditions cela marche sur un groupe quelconque.

Le troisième est l'exponentielle, en algèbre pure. Une dérivation est une
variation infinitésimale ; son exponentielle en fait une transformation
finie. En caractéristique zéro, et quand la dérivation est nilpotente, la série
s'arrête d'elle-même et tout devient algèbre. Les pages montrent que les
groupes à un paramètre qui agissent sur un module sont exactement les
exponentielles d'un opérateur, que les automorphismes d'algèbre qui
« diffèrent peu » de l'identité sont les exponentielles de dérivations, et
transportent cela aux schémas. Plusieurs énoncés y demandent une hypothèse
que la page ne met pas, et deux sont faux tels qu'écrits ; on donne ce qui
est vrai.

Le quatrième est l'orientation. Qu'est-ce qu'orienter un espace ? Pour une
droite, choisir un sens ; pour un plan, un sens de rotation ; en général,
choisir l'une de deux classes. Plutôt que de le définir, Grothendieck écrit
les règles que toute notion d'orientation doit suivre --- chaque espace en a
exactement deux, et un demi-espace oriente l'hyperplan qui le borde, comme
dans la formule de Stokes --- et montre qu'il n'y a qu'une théorie qui les
suive : celle du signe du déterminant. La dernière page commence le même
travail pour les déterminants eux-mêmes, sur un anneau quelconque.

Ce que ces pages montrent de sa manière tient en deux gestes. Le premier est
d'écrire les règles auxquelles une notion doit obéir, puis de prouver que ces
règles la déterminent : c'est la théorie de l'orientation. Le second est de
laisser la question ouverte sur la page, marquée comme telle, plutôt que de
l'enjamber ; et, quand une démonstration ne passe pas, de reprendre le texte à
zéro plutôt que de la rapiécer.

Les noms sous lesquels chercher la suite sont ceux de théorie des jeux
combinatoires et de jeux sur les graphes (théorème de Zermelo, attracteur,
détermination), de noyau d'un graphe orienté, d'inégalités de réarrangement
(Hardy--Littlewood, Riesz--Sobolev) et de majorisation, de correspondance
exponentielle entre dérivations nilpotentes et automorphismes unipotents, de
torseur d'orientation et de droite déterminant.

\keywords{combinatorial game theory, positional game, Zermelo's theorem, backward induction, reachability game, attractor, positional determinacy, Gale–Stewart theorem, kernel of a digraph, ordinal rank, rearrangement inequality, weak majorization, Riesz–Sobolev inequality, symmetric decreasing rearrangement, Pollard's theorem, exponential map, nilpotent derivation, unipotent automorphism, formal additive group, locally nilpotent derivation, differential operator, orientation torsor, determinant line, Koszul sign rule, graded determinant, Picard category}
\end{resume}

\pagerange{1}{55}
\section*{Le fil du dossier, et les conventions}

\emph{Les pièces.} Le dossier n'a pas de pagination de sa main ; on suit
celle de l'archive.

\begin{enumerate}
  \item \emph{Pages 1 à 13}, crayon : un exposé numéroté sur les jeux de
  position, §1 à §5, et le seul titre d'un §6 (section I). C'est la
  \emph{première mouture}.
  \item \emph{Pages 14 à 16} : calculs de brouillon sur des parties finies
  d'un groupe commutatif, la page 16 à l'encre bleue ; \emph{pages 18 à 21},
  encre noire puis bleue : une suite numérotée sur le réarrangement des
  fonctions positives sur un groupe (section IV). La page 17 est un tableau
  administratif dactylographié, sans rapport, que la transcription ne reprend
  pas.
  \item \emph{Pages 23 à 33}, crayon, titrées « Jeux de position » : la
  \emph{seconde mouture} de la théorie des jeux (section II).
  \item \emph{Pages 35 à 45}, crayon, titrées « Exponentielles (en car.\ 0) »
  (section V).
  \item \emph{Page 47}, crayon, « Jeux à $N$ joueurs » : une \emph{troisième
  mouture}, réduite aux données (section III).
  \item \emph{Pages 49, 51 et 53}, encre bleue : une théorie axiomatique de
  l'orientation (section VI). Elles alternent avec trois exemplaires d'une
  même page dactylographiée, pages 50, 52 et 54, intitulée « XVIII Les
  épousailles (2) » --- une liste de couples de mots, « lumière et ombre »,
  corrigée à la main, sans mathématiques, que la transcription ne reprend
  pas.\footnote{La page 52, seule des trois, porte en tête « 82 », écrit à la
  main, d'une main que le fac-similé ne permet pas d'attribuer ; comme le
  reste du tapuscrit, ce chiffre n'est pas transcrit. Le thème --- des couples
  de contraires --- est celui que Grothendieck développera sous les noms de
  yin et de yang dans \emph{Récoltes et semailles} ; le rapprochement est le
  nôtre, et rien ne dit que ce tapuscrit en soit un état.}
  \item \emph{Page 55}, crayon : les axiomes d'une théorie des déterminants
  gradués (section VII).
\end{enumerate}
Les pages 22, 34, 46 et 48 sont blanches.

\emph{Le fil.} Il n'y en a pas un, mais cinq. Les jeux (sections I à III)
forment le seul sujet repris : trois fois, et les trois rédactions sont des
\emph{moutures} distinctes, qu'on ne fond pas. La première (pages 1 à 13)
démontre le théorème central par un argument qui ne vaut qu'en l'absence de
parties infinies, s'essaie au cas général sur une page qu'il barre, et finit
sur une question. La seconde (pages 23 à 33) reprend tout, démontre le
théorème en général et va jusqu'au jeu alterné. La troisième (page 47) change
de cadre --- une position initiale, $N$ joueurs qui jouent dans un ordre
circulaire --- et s'arrête aux définitions. L'ordre dans lequel les trois ont
été écrites n'est pas établi ; celui des pages est celui qu'on suit. Le
réarrangement (section IV), l'exponentielle (section V), l'orientation et le
déterminant (sections VI et VII) sont des sujets à part ; les deux derniers se
touchent, par la règle des signes de Koszul, et ce rapprochement est dit à
la section VII.

\emph{Notations.} Chaque partie garde les lettres de ses pages, avec les
exceptions suivantes, signalées en note au point où elles interviennent. Dans
les deux premières moutures, l'application qui donne le joueur qui a
l'initiative est notée $\alpha$, comme dans la seconde (la première écrit
$i$) ; et $\overline{C}(j)$ désigne toujours l'ensemble des positions
\emph{non terminales} où l'initiative est à un autre joueur que $j$, comme
dans la seconde (la première y inclut les positions terminales). Pour le
réarrangement, la suite décroissante des valeurs d'une fonction $f$ est notée
$f^{\downarrow}$, et l'élément neutre du groupe $e$. Pour l'exponentielle, la
droite formelle est notée $\widehat{\mathbb{G}}_a$.

\emph{Collisions de lettres.} Elles sont nombreuses et on ne les corrige pas,
mais il faut les avoir en tête. $G$ est l'ensemble des positions gagnantes
(pages 31 à 33), $G(j)$ celui des positions gagnantes pour $j$, et $G$ un
groupe aux pages 14 à 21. $N$ est l'ensemble des positions nulles (pages 29
à 32), le nombre de joueurs (page 47), et $N(A)$ l'ensemble des éléments
nilpotents d'un anneau (pages 35 à 45). $P$ est l'ensemble des positions
perdantes ; $\mathfrak{P}$, l'ensemble des parties. Surtout, la théorie des
jeux combinatoires d'aujourd'hui dit \emph{N-position} (« next ») ce que la
page appelle \emph{gagnante}, \emph{P-position} (« previous ») ce qu'elle
appelle \emph{perdante}, et \emph{D-position} (« draw ») ce qu'elle appelle
\emph{nulle} : le $N$ de la page n'est pas celui de la littérature.

\emph{Ce que le dossier annonce sans l'établir.} Le §6 de la première
mouture (page 13), dont il n'y a que le titre ; le « Th 3 » de la page 12,
dans la généralité où la marge le place ; la fin de la page 21 ; l'usage des
axiomes de la page 16 ; la question du NB de la page 43 ; l'indice
$I(\omega, \omega')$ de la page 53 ; le théorème que préparent les axiomes de
la page 55 ; une condition de gain pour les jeux à $N$ joueurs de la page 47.

\pagerange{1}{13}
\section*{I. Jeux de position, première mouture (pages 1 à 13)}

\pagerange{1}{1}
\subsection*{Les données (page 1)}

Un \emph{jeu de position} est la donnée :
\begin{itemize}
  \item d'un ensemble $C$ de \emph{positions} (ou configurations) ;
  \item d'une relation $R \subset C \times C$ ; on écrit $y \in R(x)$ pour
  $(x, y) \in R$, et l'on dit que $y$ est un \emph{successeur} de $x$. Les
  positions sans successeur sont dites \emph{terminales} :
  $C_0 = \{x \in C \mid R(x) = \emptyset\}$ ;
  \item d'un ensemble $J$ de \emph{joueurs} et d'une application
  $\alpha : C \smallsetminus C_0 \to J$, où $\alpha(x)$ est le joueur qui a
  l'\emph{initiative} dans la position $x$ ; on pose
  $C(j) = \alpha^{-1}(j)$, d'où
  $C \smallsetminus C_0 = \bigsqcup_{j \in J} C(j)$, et
  $\overline{C}(j) = (C \smallsetminus C_0) \smallsetminus C(j)$ ;
  \item pour chaque joueur $j$, d'une partie $G_0(j) \subset C_0$, les
  positions terminales \emph{gagnantes pour $j$}.
\end{itemize}
Rien n'oblige les $G_0(j)$ à être disjoints ; une position terminale peut
n'être gagnante pour personne.\footnote{La page note l'initiative $i$ ; on
écrit $\alpha$, la lettre de la seconde mouture. L'application était d'abord
définie sur $C$ tout entier : la restriction à $C \smallsetminus C_0$ est
ajoutée à l'encre bleue, et une remarque qui le justifiait est barrée. La page
pose $\overline{C}(j) = C \smallsetminus C(j)$, qui contient $C_0$, et écrit
ensuite $\overline{C}(j) \smallsetminus \overline{C}_0(j)$ pour les positions
non terminales ; on adopte d'emblée la convention de la page 23. Une note
marginale, dont la négation est d'une lecture incertaine, précise qu'« on ne
suppose pas $G_0(j) \subset C_0(j)$ » ; $C_0(j)$ n'est défini nulle part, et
la remarque n'a de sens que pour la version où $\alpha$ était définie sur
$C$ : une position terminale gagnante pour $j$ n'a pas à être une position
« où $j$ a la main ».}

\pagerange{1}{2}
\subsection*{Parties, et positions de longueur bornée (pages 1 et 2)}

Une \emph{partie} à $n$ coups est une suite $(x_0, \ldots, x_n)$ avec
$x_{i+1} \in R(x_i)$ ; elle ne dépend que de $(C, R)$. Elle est
\emph{terminée} si $x_n \in C_0$. On parle aussi de parties infinies, et une
partie est \emph{circulaire} si $x_n = x_0$ avec $n \geqslant 1$. S'il existe
une partie circulaire, il existe une partie infinie ; la réciproque vaut si
$C$ est fini, puisqu'une partie infinie repasse alors par une même position.

On définit par récurrence
\[
C_{-1} = \emptyset,
\qquad
C_i = \{x \in C \mid R(x) \subset C_{i-1}\} \quad (i \geqslant 0),
\]
de sorte que $C_0$ est bien l'ensemble des positions terminales. La suite
$(C_i)$ est croissante, et $C_i$ est l'ensemble des positions dont toutes les
parties issues sont de longueur au plus $i$. La réunion
$C_\infty = \bigcup_i C_i$ est donc l'ensemble des positions d'où les
longueurs des parties sont \emph{bornées}. Il revient au même de demander
$C = C_\infty$ ou l'existence d'une fonction $\ell : C \to \mathbb{N}$ qui
décroît strictement le long des coups ($\ell(y) < \ell(x)$ pour
$y \in R(x)$) ; on a alors $x \in C_{\ell(x)}$.\footnote{La page 2 écrit « pour
tout $x$, l'une des longueurs des parties partant de $x$ est finie » ; ce qui
est équivalent à $C = C_\infty$ est que ces longueurs soient bornées, comme le
dit le corollaire qui précède.}

Ces conditions interdisent les parties infinies. Si $C$ est fini, de cardinal
$N$, la réciproque est vraie : sans partie infinie, il n'y a pas de partie
circulaire, une partie ne repasse jamais par la même position, sa longueur est
au plus $N - 1$, et $C = C_{N-1}$.\footnote{Les mots « la réciproque est
vraie » sont d'une lecture incertaine, sur une ligne surchargée ; c'est bien
ce que dit la suite. Pour $C$ infini, l'absence de partie infinie
n'entraîne pas $C = C_\infty$ : une position dont les successeurs $y_1, y_2,
\ldots$ commencent des chaînes forcées de longueurs $1, 2, \ldots$ n'a que des
parties finies, mais de longueurs non bornées (exemple nôtre). La bonne
généralisation est transfinie : $C_\alpha = \{x \mid R(x) \subset
\bigcup_{\beta < \alpha} C_\beta\}$ pour tout ordinal $\alpha$, et l'absence
de partie infinie équivaut à $C = \bigcup_\alpha C_\alpha$ (avec l'axiome des
choix dépendants) ; le plus petit $\alpha$ tel que $x \in C_\alpha$ est le
\emph{rang} ordinal de $x$. Ce point commande la question de la page 12.}

\pagerange{2}{2}
\subsection*{Stratégies et positions gagnantes (page 2)}

Une \emph{stratégie} pour le joueur $j$ est une application
$\Sigma : C(j) \to \mathfrak{P}(C)$ telle que
$\emptyset \neq \Sigma(x) \subset R(x)$ pour tout $x \in C(j)$ : en chaque
position où il a la main, $j$ s'autorise un ensemble non vide de coups. Une
partie est \emph{compatible} avec $\Sigma$ si chacun de ses coups joués depuis
une position de $C(j)$ est dans $\Sigma$ ; même définition pour les parties
infinies. La stratégie $\Sigma$ est \emph{gagnante pour $x$} si
\begin{itemize}
  \item[a)] il n'y a pas de partie infinie issue de $x$ et compatible avec
  $\Sigma$ ;
  \item[b)] toute partie terminée issue de $x$ et compatible avec $\Sigma$
  s'achève dans $G_0(j)$.
\end{itemize}
La position $x$ est \emph{gagnante pour $j$} s'il existe une stratégie pour
$j$ gagnante pour $x$ ; on note $G(j)$ l'ensemble de ces positions. On a
$G(j) \cap C_0 = G_0(j)$.

Deux traits de cette définition sont modernes avant la lettre. Une stratégie
ne regarde que la position présente, non l'histoire de la partie : c'est ce
qu'on appelle aujourd'hui une stratégie \emph{positionnelle}, ou sans
mémoire. Et elle autorise un ensemble de coups plutôt qu'un seul, ce qu'on
appelle parfois une multi-stratégie. Le second trait dispense de choisir un
coup en chaque position ; le premier se révélera sans perte de généralité
(section II).

\pagerange{3}{6}
\subsection*{Théorème 1 : l'équation de point fixe (pages 3 à 6)}

\textbf{Théorème 1.} Soit $j \in J$.
\begin{itemize}
  \item[1°)] Pour $x \in C(j)$ : $x \in G(j)$ si et seulement si
  $R(x) \cap G(j) \neq \emptyset$.
  \item[2°)] Pour $x \in \overline{C}(j)$ : $x \in G(j)$ si et seulement si
  $R(x) \subset G(j)$.
\end{itemize}

Avec $G(j) \cap C_0 = G_0(j)$, cela dit que $G(j)$ est un point fixe de
l'opérateur, croissant pour l'inclusion,
\[
F_j(X) = G_0(j)
\cup \{x \in C(j) \mid R(x) \cap X \neq \emptyset\}
\cup \{x \in \overline{C}(j) \mid R(x) \subset X\},
\]
qu'on introduit pour la commodité de ce qui suit.\footnote{L'opérateur $F_j$
n'est pas sur la page ; c'est la forme que prend aujourd'hui le calcul des
régions gagnantes d'un jeu sur un graphe.}

Le théorème est vrai tel quel, sans hypothèse. La démonstration de la page ne
l'établit pourtant qu'en partie.\footnote{La marge de la page 3 porte « S'il
n'y a pas de partie infinie », la négation étant d'une lecture incertaine :
c'est l'hypothèse sous laquelle la démonstration de 2°, $\Leftarrow$, est
correcte.} Elle repose sur un lemme : si $\Sigma$ est gagnante pour $x$, elle
l'est pour toute position d'une partie issue de $x$ et compatible avec
$\Sigma$ ; en particulier $\Sigma(x) \subset G(j)$ quand $x \in C(j)$.

\emph{1°.} Le sens direct est le lemme. Pour la réciproque, soit
$y \in R(x) \cap G(j)$ et $\Sigma$ gagnante pour $y$ ; on la modifie en
posant $\Sigma'(x) = \{y\}$. Si une partie issue de $y$ et compatible avec
$\Sigma$ passe par $x$, alors $x$ est déjà gagnante pour $\Sigma$ par le
lemme ; sinon, une partie issue de $x$ et compatible avec $\Sigma'$ ne revient
jamais en $x$, suit $\Sigma$ après son premier coup, et est
gagnée.\footnote{C'est la note écrite en oblique dans la marge de la page 4 :
« Montrons d'abord qu'il n'y a pas de partie compatible avec $\Sigma$
commençant avec $y$ et terminant par $x$ », les derniers mots d'une lecture
incertaine. Plusieurs mots de l'argument sont illisibles ; l'articulation
donnée ici est celle que la page 26 rédige en clair.}

\emph{2°.} Le sens direct est immédiat : une stratégie gagnante pour $x$ l'est
pour chaque $y \in R(x)$, puisque $x$ n'est pas une position de $j$. Pour la
réciproque, la page 5 propose la stratégie
\[
\Sigma(z) = R(z) \cap G(j) \ \ (z \in C(j) \cap G(j)),
\qquad
\Sigma(z) = R(z) \ \ \text{sinon},
\]
et montre par récurrence que toute partie issue de $x$ et compatible avec
elle reste dans $G(j)$ ; une telle partie, si elle se termine, se termine
donc dans $G(j) \cap C_0 = G_0(j)$. Ce raisonnement ne montre pas qu'elle se
termine, et ne suffit donc qu'en l'absence de parties infinies. Rester dans
les positions gagnantes n'est pas gagner.\footnote{Exemple nôtre : deux
positions $a, b \in C(j)$ avec $R(a) = \{b, t\}$, $R(b) = \{a\}$, et
$t \in G_0(j)$. Les deux sont gagnantes ($a$ en jouant $t$, $b$ en jouant
$a$ puis $t$), mais la stratégie « rester dans $G(j)$ », qui autorise
$a \to b$, laisse jouer $a, b, a, b, \ldots$ indéfiniment. Le mot
« convenable » qui qualifie la stratégie, et la note marginale de la page 5,
lue « On suppose qu'il y a une partie infinie », sont d'une lecture
incertaine ; on attendrait plutôt « qu'il n'y a pas ».}

La page 6, traversée d'un trait oblique qui l'annule, tente le cas général :
on munit $R(x)$ d'un bon ordre, on choisit pour chaque $y \in R(x)$ une
stratégie $\Sigma_y$ gagnante pour $y$, et, en une position $z$ atteinte, on
suit $\Sigma_{y(z)}$ pour le plus petit $y(z)$ d'où $z$ est atteint ; le
long d'une partie, les $y(x_i)$ décroissent, donc stationnent. La page
s'interrompt là. Telle qu'écrite, l'accessibilité est prise au sens de
parties quelconques, et la queue de partie qui suit $\Sigma_y$ après la
stationnarité n'a alors aucune raison d'être gagnante pour $\Sigma_y$. La
seconde mouture reprend exactement cet argument, en prenant l'accessibilité
par des parties compatibles avec $\Sigma_y$, et il passe (pages 26 à 28,
section II).

\pagerange{5}{8}
\subsection*{La plus grande stratégie gagnante (pages 5, 7 et 8)}

Disons qu'une partie $\Gamma \subset G(j)$ est \emph{stable} si
$R(z) \cap G(j) \subset \Gamma$ pour $z \in \Gamma \cap C(j)$ et
$R(z) \subset \Gamma$ pour $z \in \Gamma \cap \overline{C}(j)$ ; par le
théorème 1, $G(j)$ lui-même est stable. On ordonne les stratégies par
inclusion : $\Sigma' \subset \Sigma$ si $\Sigma'(z) \subset \Sigma(z)$ pour
tout $z$.

\textbf{Corollaire.} Supposons qu'il n'y ait pas de partie infinie, et soit
$\Gamma \subset G(j)$ stable. La stratégie
\[
\Sigma_\Gamma(z) = R(z) \cap G(j) \ \ (z \in \Gamma \cap C(j)),
\qquad
\Sigma_\Gamma(z) = R(z) \ \ (z \in C(j) \smallsetminus \Gamma)
\]
est gagnante pour tout $x \in \Gamma$, et elle est la plus grande des
stratégies gagnantes pour tout $x \in \Gamma$.

C'est une stratégie par le 1° du théorème 1. Une partie issue de
$x \in \Gamma$ et compatible avec elle reste dans $\Gamma$ par stabilité ;
elle est finie par hypothèse, et se termine dans
$\Gamma \cap C_0 \subset G_0(j)$. Enfin, si $\Sigma'$ est gagnante pour tout
point de $\Gamma$, le lemme donne $\Sigma'(z) \subset G(j)$ pour
$z \in \Gamma \cap C(j)$, donc $\Sigma' \subset \Sigma_\Gamma$ (pages 7 et
8).\footnote{Le mot « stable » est d'une lecture incertaine et n'est pas
défini ; la définition est la nôtre. Elle est lue dans les conditions
biffées de la page 7, qui cherchaient à restreindre la stratégie aux
positions accessibles depuis $\Gamma$ ; sans elle, la démonstration de la page
7, qui applique la définition de $\Sigma$ en des positions $x_i \in C(j)$
comme si elles étaient dans $\Gamma$, ne passe pas. La page écrit aussi
$R(z)$ pour $R(x_i)$. La démonstration, comme celle du 2° du théorème 1, ne
vérifie que la condition b) de gain ; l'hypothèse de l'absence de parties
infinies, qui n'est pas écrite à cet endroit, est nécessaire.}

Sans cette hypothèse, il n'y a pas en général de plus grande stratégie
gagnante.\footnote{Exemple nôtre : $a, b \in C(j)$, $R(a) = \{b, t\}$,
$R(b) = \{a, t'\}$, $t, t' \in G_0(j)$. Les stratégies
$(\Sigma(a), \Sigma(b)) = (\{b, t\}, \{t'\})$ et $(\{t\}, \{a, t'\})$ sont
gagnantes pour $a$ et pour $b$ ; toute stratégie qui les contient autorise la
boucle $a, b, a, \ldots$. La notion est celle de stratégie \emph{la plus
permissive} (Bernet, Janin et Walukiewicz, 2002) : elle existe pour les
objectifs de sûreté, non en général pour les objectifs d'atteinte comme
celui-ci.}

\pagerange{8}{11}
\subsection*{Positions acceptables ; théorème 2 (pages 8 à 11)}

Pour $J' \subset J$, une \emph{$J'$-stratégie} se définit comme une stratégie
pour un joueur, en remplaçant $C(j)$ par $C(J') = \bigcup_{j \in J'} C(j)$ :
c'est la stratégie d'une coalition. Fixons $j' \in J$ et
$J' = J \smallsetminus \{j'\}$. Une $J'$-stratégie $\Sigma'$ est
\emph{acceptable pour $x$ contre $j'$} si toute partie terminée issue de $x$
et compatible avec $\Sigma'$ s'achève hors de $G_0(j')$ ; et $x$ est
\emph{acceptable contre $j'$} s'il existe une telle $\Sigma'$. Les parties
infinies ne comptent pas contre la coalition : être acceptable, c'est
empêcher $j'$ de gagner.

\textbf{Théorème 2.} Pour que $x \notin G(j')$, il faut et il suffit que $x$
soit acceptable contre $j'$.

\emph{Suffisance} (page 9). Si $x \in G(j')$ par une stratégie $\Sigma$ et
que $\Sigma'$ est acceptable pour $x$, une partie maximale issue de $x$ et
compatible avec les deux est finie, puisque $\Sigma$ est gagnante, et
s'achève à la fois dans $G_0(j')$ et hors de $G_0(j')$.\footnote{La page écrit
$G(j)$ là où l'argument demande $G_0(j')$ ; $j$ et le prime sont peu
distincts sur cette ligne. L'existence d'une partie maximale compatible avec
les deux stratégies, qui est implicite ici, est le lemme de la page 24 ; elle
repose sur l'axiome des choix dépendants.}

\emph{Nécessité} (pages 10 et 11). On pose, pour $z \in C(J')$,
\[
\Sigma'(z) = R(z) \ \ (z \in G(j')),
\qquad
\Sigma'(z) = R(z) \smallsetminus G(j') \ \ (z \notin G(j')).
\]
C'est une $J'$-stratégie par le 2° du théorème 1, et une partie issue d'une
position $x \notin G(j')$ et compatible avec elle reste hors de $G(j')$ : aux
positions de $j'$ par le 1°, aux autres par construction. Elle ne peut donc
s'achever dans $G_0(j')$. Cette stratégie est acceptable pour \emph{toutes}
les positions non gagnantes pour $j'$ à la fois.

Le théorème 2 vaut sans hypothèse, dès lors que le théorème 1 vaut en
général. Grothendieck observe alors qu'il a « fait deux fois le même
raisonnement », et réduit le théorème au jeu auxiliaire à deux joueurs
$J(j')$ : $j'$ contre la coalition, notée $\bar\jmath$, avec les mêmes
positions et les mêmes coups, $C(\bar\jmath) = C(J')$,
$G_0(\bar\jmath) = C_0 \smallsetminus G_0(j')$. Les $J'$-stratégies du jeu
initial sont les stratégies de $\bar\jmath$ dans $J(j')$, et, \emph{s'il n'y
a pas de partie infinie}, être acceptable contre $j'$ revient à être gagnant
pour $\bar\jmath$.\footnote{L'hypothèse est dans la marge de la page 11, et
elle est nécessaire : une partie infinie compte pour la coalition dans
« acceptable », non dans « gagnante ». Ce décalage est celui qui sépare, dans
le vocabulaire d'aujourd'hui, un objectif d'\emph{atteinte} (ouvert) d'un
objectif de \emph{sûreté} (fermé).} On en tire, sans parties infinies,
qu'en tout jeu à deux joueurs où chaque position terminale est gagnante pour
exactement l'un des deux, toute position est gagnante pour exactement l'un
des deux. C'est le théorème de Zermelo, qui est ce que la page écrit, en
bas, par « gagnante pour $j$ $\Rightarrow$ non gagnante pour $j'$ » et
l'implication inverse.\footnote{Zermelo (1913), pour les échecs ; König
(1927) et Kalmár (1928) pour les jeux où toute partie est finie. Le théorème 2
lui-même est un cas du théorème de Gale et Stewart (1953) : un jeu dont
l'objectif est ouvert est déterminé ; ici avec des stratégies positionnelles
des deux côtés. Rien sur la page ne cite ces travaux.}

\pagerange{12}{12}
\subsection*{« ? Th 3 » : la construction récurrente (page 12)}

La page 12 cherche à construire $G(j)$ comme la page 1 construisait les
$C_i$. On pose $G_{-1}(j) = \emptyset$ et
\[
G_i(j) = F_j\bigl(G_{i-1}(j)\bigr) \qquad (i \geqslant 0),
\]
c'est-à-dire : $G_i(j)$ est l'ensemble des positions d'où $j$ peut forcer le
gain en au plus $i$ coups. La suite est croissante, contenue dans $G(j)$ par
le théorème 1, et l'on pose $G_\omega(j) = \bigcup_i G_i(j) \subset G(j)$.
La page énonce alors, sous un point d'interrogation :

\begin{quote}
\textbf{? Th 3.} $G_\omega(j) = G(j)$.
\end{quote}

et démontre, par récurrence sur $i$ au moyen du théorème 1, que
$G_\omega(j) \cap C_i = G(j) \cap C_i$ pour tout $i$.\footnote{La page note
$G_\infty(j)$ ; on écrit $G_\omega(j)$ pour la distinguer de l'itération
transfinie. Une première tentative de démonstration, qui partait d'une
stratégie gagnante et d'une partie maximale, est encadrée, hachurée et
abandonnée.} Voici ce qui est vrai.

\begin{itemize}
  \item[a)] Si $C = \bigcup_i C_i$ --- les longueurs des parties issues de
  chaque position sont bornées ---, alors $G_\omega(j) = G(j)$ : c'est la
  démonstration de la page. Le cas d'un $C$ fini sans partie circulaire est
  compris.
  \item[b)] Si $C$ est fini, même avec des parties circulaires, alors
  $G_\omega(j) = G(j)$ : la suite croissante $G_i(j)$ stationne, et sa limite
  est le plus petit point fixe de $F_j$, qui est $G(j)$ (point d) ci-dessous).
  \item[c)] Plus généralement, $G_\omega(j) = G(j)$ dès que chaque position
  de $\overline{C}(j)$ n'a qu'un nombre fini de successeurs.
  \item[d)] En général, $G_\omega(j) \neq G(j)$, même quand il n'y a aucune
  partie infinie. Il faut prolonger la récurrence au-delà des entiers :
  \[
  G_\alpha(j) = F_j\Bigl(\,\bigcup_{\beta < \alpha} G_\beta(j)\Bigr)
  \quad (\alpha \text{ ordinal}),
  \qquad
  G(j) = \bigcup_\alpha G_\alpha(j),
  \]
  et $G(j)$ est le plus petit point fixe de $F_j$.
\end{itemize}

Les points b) à d) sont les nôtres.\footnote{La marge de la page 12 énonce le
théorème « si $C$ fini, ou s'il n'existe pas de partie infinie, plus
généralement si $C = \bigcup C_i$ ». La démonstration couvre le dernier cas,
mais l'absence de parties infinies n'en est pas un cas particulier quand $C$
est infini (note de la page 2). Contre-exemple (nôtre) : une position
$x \in \overline{C}(j)$ dont les successeurs $y_1, y_2, \ldots$ commencent des
chaînes forcées de longueurs $1, 2, \ldots$ aboutissant dans $G_0(j)$. Toute
partie est finie et gagnée par $j$, donc $x \in G(j)$ ; mais
$y_n \in G_n(j) \smallsetminus G_{n-1}(j)$, de sorte que
$R(x) \not\subset G_i(j)$ pour tout $i$ et $x \notin G_\omega(j)$ ; on a
$x \in G_{\omega + 1}(j)$.} Pour d), on voit que $G_\alpha(j) \subset G(j)$
par récurrence transfinie à l'aide du théorème 1 ; inversement, si $\Sigma$
est gagnante pour $x$, l'arbre des parties issues de $x$ et compatibles avec
$\Sigma$ n'a pas de branche infinie, et une récurrence sur son rang montre
que $x$ appartient à un $G_\alpha(j)$. La suite transfinie stationne à un
ordinal de cardinal au plus celui de $C$. Pour c), l'hypothèse rend $F_j$
compatible aux réunions croissantes dénombrables, et le point fixe est atteint
en $\omega$ étapes.\footnote{Dans le vocabulaire des jeux sur les graphes,
$G(j)$ est l'\emph{attracteur} de $G_0(j)$ pour le joueur $j$, et les
$G_\alpha(j)$ en sont les étapes ; sur un graphe infini, l'attracteur se
calcule par récurrence transfinie. Le point d) donne aussi une démonstration
du théorème 1 qui n'a besoin ni de bon ordre ni de choisir des stratégies
(seul l'axiome des choix dépendants intervient) : la stratégie du rang, qui en une position
$z \in C(j) \cap G(j)$ autorise les coups vers les positions de rang plus
petit, est gagnante, parce que le rang décroît strictement le long des
parties.}

\pagerange{13}{13}
\subsection*{Un paragraphe annoncé (page 13)}

La page 13 ne porte que le titre d'un §6, « Élimination des parties
circulaires (en $C$ fini) ». Le paragraphe n'est écrit nulle part dans le
dossier, et l'on n'en reconstitue pas le contenu. On note seulement qu'en $C$
fini, le point b) ci-dessus rend la question du « Th 3 » sans objet.

\pagerange{23}{33}
\section*{II. Jeux de position, seconde mouture (pages 23 à 33)}

\pagerange{23}{24}
\subsection*{Données, parties, stratégies (pages 23 et 24)}

Les données sont celles de la première mouture, écrites cette fois avec la
convention $\overline{C}(j) = (C \smallsetminus C_0) \smallsetminus C(j)$
dès le départ.\footnote{La page 23 écrit la relation avec un $R$ rond ; les
pages suivantes ont un $R$ droit.} Une \emph{partie} est une suite finie ou
infinie de positions dont deux termes consécutifs vérifient
$x_{i+1} \in R(x_i)$ ; les parties \emph{terminées}, ou non prolongeables,
sont les parties infinies et les parties finies dont la dernière position est
terminale. Une marge note que les parties se composent de façon associative,
et parle de « catégorie des positions » : c'est la catégorie libre sur le
graphe $(C, R)$, dont les parties finies sont les flèches. Une partie est
\emph{gagnée pour $j$} si elle est finie, terminée, et s'achève dans
$G_0(j)$ ; une marge observe qu'une partie composée $XY$ est gagnée si et
seulement si $Y$ l'est.

Une stratégie pour $j$ est, comme en I, une application
$\Sigma : C(j) \to \mathfrak{P}(C)$ avec
$\emptyset \neq \Sigma(x) \subset R(x)$ ; on la prolonge par
$\Sigma(x) = R(x)$ sur $\overline{C}(j)$. Les parties compatibles avec
$\Sigma$ --- les « $\Sigma$-parties » --- sont stables par composition et
forment une sous-catégorie.

\textbf{Lemme} (page 24). Soient $J' \subset J$ et, pour chaque $j \in J'$,
une stratégie $\Sigma_j$ pour $j$. De toute position part une partie
terminée compatible avec toutes les $\Sigma_j$.\footnote{On la construit coup
par coup : c'est l'axiome des choix dépendants, que la page ne mentionne
pas.}

La stratégie $\Sigma$ est \emph{gagnante pour $x_0$} si toute partie
terminée issue de $x_0$ et compatible avec $\Sigma$ est gagnée pour $j$ :
c'est la définition de la section I, les deux conditions a) et b) réunies en
une. On note encore $G(j)$ l'ensemble des positions gagnantes pour $j$, et
$G(j) \cap C_0 = G_0(j)$. Une proposition entre crochets en tire, par le
lemme, que si $\bigcap_{j \in J'} G_0(j) = \emptyset$ (avec $J'$ non vide),
aucune position n'est gagnante pour tous les $j \in J'$ à la fois.

\pagerange{25}{28}
\subsection*{Le théorème de point fixe, démontré (pages 25 à 28)}

\textbf{Proposition} (page 25). Si $\Sigma$ est gagnante pour $x_0$, toute
position d'une $\Sigma$-partie issue de $x_0$ est gagnante pour $\Sigma$.

\textbf{Théorème} (page 25). On a
\[
G(j) \cap C(j) = \{x \in C(j) \mid R(x) \cap G(j) \neq \emptyset\},
\qquad
G(j) \cap \overline{C}(j) = \{x \in \overline{C}(j) \mid R(x) \subset G(j)\},
\]
et $G(j) \cap C_0 = G_0(j)$ ; autrement dit, $G(j) = F_j(G(j))$.

C'est le théorème 1 de la page 3, et la démonstration en est cette fois
complète. Le point a), $\supset$, est celui de la section I : si
$x_1 \in R(x_0) \cap G(j)$ par une stratégie $\Sigma_1$, et si $x_0$ n'est
pas déjà gagnante pour $\Sigma_1$, la stratégie $\Sigma_0$ égale à
$\Sigma_1$ sauf en $x_0$, où $\Sigma_0(x_0) = \{x_1\}$, est gagnante : une
partie jouée suivant $\Sigma_0$ ne revient jamais en $x_0$, car le premier
retour ferait de $x_0$ une position gagnante pour $\Sigma_1$.

Le point b), $\supset$, est le cœur. Soit $x_0 \in \overline{C}(j)$ avec
$R(x_0) \subset G(j)$ ; on choisit pour chaque $x_1 \in R(x_0)$ une stratégie
$\Sigma_{x_1}$ gagnante pour $x_1$, et un bon ordre sur $R(x_0)$. Pour
$x \in C$, soit $E(x)$ l'ensemble des $x_1 \in R(x_0)$ tels qu'il existe une
$\Sigma_{x_1}$-partie de $x_1$ à $x$, soit
$C' = \{x \mid E(x) \neq \emptyset\}$, et, pour $x \in C'$, soit $\varepsilon(x)$
le plus petit élément de $E(x)$. On pose
\[
\Sigma_0(x) = \Sigma_{\varepsilon(x)}(x) \ \ (x \in C(j) \cap C'),
\qquad
\Sigma_0(x) = R(x) \ \ \text{sinon}.
\]
Un lemme (page 27) dit que tout $x \in C'$ est gagnant pour
$\Sigma_{\varepsilon(x)}$, donc $C' \subset G(j)$, et que pour
$y \in \Sigma_0(x)$ on a $\varepsilon(x) \in E(y)$, donc $y \in C'$ et
$\varepsilon(y) \leqslant \varepsilon(x)$. Le long d'une partie issue de
$x_0$ et compatible avec $\Sigma_0$, les $\varepsilon(x_i)$, $i \geqslant 1$,
forment alors une suite décroissante dans un ensemble bien ordonné : elle
stationne en une valeur $\xi$ à partir d'un rang $n$, et la queue
$(x_n, x_{n+1}, \ldots)$ est une $\Sigma_\xi$-partie issue d'une position
gagnante pour $\Sigma_\xi$ ; elle est donc finie. Comme toutes les $x_i$ sont
dans $C' \subset G(j)$, la partie s'achève dans $G(j) \cap C_0 = G_0(j)$.

C'est très exactement l'argument que la page 6 avait ébauché puis barré ; ce
qui le fait passer est de mesurer l'accessibilité par des parties compatibles
avec $\Sigma_{x_1}$, et non par des parties quelconques.\footnote{L'indice du
point de départ, dans « de $x_1$ à $x$ », est surchargé ; $x_1$ est la
lecture que demande la suite. Le choix des $\Sigma_{x_1}$ et du bon ordre
utilise l'axiome du choix ; la stratégie du rang (note de la page 12) s'en
dispense.}

\emph{Lecture moderne.} Le théorème dit que $G(j)$ est un point fixe de
$F_j$ ; il ne dit pas lequel, et les équations seules ne le caractérisent pas
quand il y a des parties infinies.\footnote{Deux positions $a, b \in C(j)$
avec $R(a) = \{b\}$, $R(b) = \{a\}$ et $G_0(j) = \emptyset$ : $X = \{a, b\}$
vérifie $F_j(X) = X$, alors que $G(j) = \emptyset$ (exemple nôtre).} Ce qui
le caractérise est d'être le \emph{plus petit} point fixe (section I, page 12,
point d)) ; quand il n'y a pas de partie infinie, le point fixe est unique,
par récurrence sur le rang. C'est ce qu'établit la page 33 dans le cas
alterné.

\pagerange{28}{29}
\subsection*{Stratégies contre $j$ ; positions admissibles (pages 28 et 29)}

Une \emph{stratégie contre $j$} est une application $\Sigma'$ définie sur
$C \smallsetminus C_0$, avec $\Sigma'(x) = R(x)$ pour $x \in C(j)$ et
$\emptyset \neq \Sigma'(x) \subset R(x)$ pour $x \in \overline{C}(j)$ :
c'est la stratégie de la coalition de tous les autres joueurs. Elle est
\emph{admissible pour $x_0$} si toute partie terminée issue de $x_0$ et
jouée suivant $\Sigma'$ n'est pas gagnée par $j$, c'est-à-dire est infinie
ou s'achève hors de $G_0(j)$.\footnote{La page écrit « jouée suivant
$x_0$ » ; on lit « suivant $\Sigma'$ ».} C'est la notion que la première
mouture appelait « acceptable ».

\textbf{Théorème} (pages 28 et 29). Soit $\Sigma'_*$ la stratégie contre $j$
définie par
\[
\Sigma'_*(x) = R(x) \ \ (x \in C(j) \cup G(j)),
\qquad
\Sigma'_*(x) = R(x) \smallsetminus G(j) \quad
(x \in \overline{C}(j) \smallsetminus G(j)).
\]
Pour $x_0 \in C$, les conditions suivantes sont équivalentes : a)
$x_0 \notin G(j)$ ; b) $x_0$ est admissible contre $j$ ; c) $x_0$ est
admissible pour $\Sigma'_*$.

C'est une stratégie par le b) du théorème de point fixe ; c)
$\Rightarrow$ b) est trivial, b) $\Rightarrow$ a) résulte du lemme de la page
24, et pour a) $\Rightarrow$ c) une partie jouée suivant $\Sigma'_*$ depuis
$x_0 \notin G(j)$ reste hors de $G(j)$.\footnote{La page énonce les conditions
« pour $x \in C$ » et les écrit pour $x_0$.} C'est le théorème 2 de la
première mouture, et il contient une remarque qui n'est pas sur la page : une
seule stratégie positionnelle de la coalition bloque $j$ depuis toutes ses
positions non gagnantes, contre \emph{toute} façon de jouer de $j$, avec ou
sans mémoire. Restreindre $j$ aux stratégies positionnelles ne lui retire donc
rien : $G(j)$ est aussi l'ensemble des positions d'où $j$ gagne avec des
stratégies qui tiennent compte de l'histoire. C'est la \emph{détermination
positionnelle} des jeux d'atteinte.

\textbf{Corollaire} (page 29). Supposons $J = \{j, j'\}$ et
$G_0(j) \cap G_0(j') = \emptyset$, et soit $N$ l'ensemble des positions
admissibles contre $j$ et contre $j'$. Alors
\[
C = G(j) \sqcup G(j') \sqcup N .
\]
Les positions de $N$ sont celles d'où aucun des deux joueurs ne peut forcer
le gain : la partie peut se prolonger indéfiniment, ou finir sur une position
qui ne donne la victoire à personne.

\pagerange{30}{32}
\subsection*{Gagnantes, perdantes, nulles : le jeu alterné (pages 30 à 32)}

Toujours à deux joueurs, notons $\bar\alpha(x)$ l'adversaire du joueur
$\alpha(x)$ qui a la main en $x \in C \smallsetminus C_0$. La page 30 pose
\[
G^* = \{x \in C \smallsetminus C_0 \mid x \in G(\alpha(x))\},
\qquad
P^* = \{x \in C \smallsetminus C_0 \mid x \in G(\bar\alpha(x))\},
\]
les positions non terminales \emph{gagnantes} (pour celui qui a la main) et
\emph{perdantes} (gagnantes pour l'autre), et appelle \emph{nulles} les
autres : $N^* = N \cap (C \smallsetminus C_0)$. On a
$G(j) \cap C(j) = G^* \cap C(j)$ et
$G(j) \cap \overline{C}(j) = P^* \cap C(j')$, de sorte que les $G(j)$ se
reconstituent à partir de $G^*$, de $P^*$ et des $G_0(j)$.

Le jeu est \emph{alterné} s'il existe $\hat\alpha : C \to J$ prolongeant
$\alpha$ avec $\hat\alpha(y) \neq \hat\alpha(x)$ pour $y \in R(x)$. Il faut et
il suffit que deux coups consécutifs vers des positions non terminales
changent de joueur, et que les prédécesseurs d'une même position terminale
appartiennent tous au même joueur ; $\hat\alpha$ est alors déterminée hors des
positions \emph{isolées}, sans prédécesseur ni successeur, où l'on peut la
choisir librement. On pose
\[
G = \{x \in C \mid x \in G(\hat\alpha(x))\},
\qquad
P = \{x \in C \mid x \in G(\bar{\hat\alpha}(x))\},
\]
qui prolongent $G^*$ et $P^*$, et l'on a $C \smallsetminus (G \cup P) = N$.
Les $G(j)$ se lisent sur $\hat\alpha$, $G$ et $P$, et pour
$x \in C \smallsetminus C_0$ :
\[
x \in G \iff R(x) \cap P \neq \emptyset,
\qquad
x \in P \iff R(x) \subset G .
\]
Une position est gagnante si l'on peut jouer vers une position perdante, et
perdante si tous les coups mènent à des positions gagnantes : c'est
l'analyse rétrograde des jeux combinatoires.\footnote{Dans le vocabulaire de
la théorie des jeux combinatoires (Berlekamp, Conway et Guy, \emph{Winning
Ways}, 1982), ce que la page appelle gagnante, perdante et nulle se dit
N-position, P-position et, pour les jeux où les parties peuvent ne pas finir
(\emph{loopy games}), D-position. Dans la convention dite normale, où celui
qui ne peut plus jouer perd, $G_0(j)$ est l'ensemble des positions terminales
où c'est à l'adversaire de $j$ de jouer, et toute position terminale est
perdante. Le $G_0$ de la page laisse cette convention libre.}

\pagerange{33}{33}
\subsection*{Caractérisation sans parties infinies (page 33)}

\textbf{Proposition} (page 33). Soit un jeu alterné à deux joueurs, et soient
$\widetilde{G}, \widetilde{P} \subset C$ tels que
\begin{itemize}
  \item[1)] pour $x \in \widetilde{G} \smallsetminus C_0$, il existe
  $y \in R(x) \cap \widetilde{P}$ ;
  \item[2)] pour $x \in \widetilde{P}$, $R(x) \subset \widetilde{G}$ ;
  \item[3)] pour $x \in \widetilde{G} \cap C_0$, $x \in G_0(\hat\alpha(x))$ ;
  \item[4)] pour $x \in \widetilde{P} \cap C_0$,
  $x \in G_0(\bar{\hat\alpha}(x))$.
\end{itemize}
S'il n'y a pas de partie infinie, alors $\widetilde{G} \subset G$ et
$\widetilde{P} \subset P$. Si de plus $\widetilde{G}$ et $\widetilde{P}$ sont
complémentaires et $G_0(j) \cap G_0(j') = \emptyset$, alors
$\widetilde{G} = G$, $\widetilde{P} = P$, et par suite $G \cup P = C$ et
$C_0 = G_0(j) \sqcup G_0(j')$.\footnote{La page énonce 1) pour tout
$x \in \widetilde{G}$. Pris à la lettre, il interdit à $\widetilde{G}$ de
contenir une position terminale, rend 3) vide, et la seconde conclusion
devient fausse dès qu'une position terminale est gagnante pour celui qui
aurait la main. On restreint 1) aux positions non terminales, ce que demande
la présence de 3).}

Pour la première assertion, chaque joueur joue, depuis ses positions de
$\widetilde{G}$, vers $\widetilde{P}$ : la partie alterne entre
$\widetilde{G}$ et $\widetilde{P}$, se termine faute de parties infinies, et
les conditions 3) et 4) disent que la position finale donne la victoire au
bon joueur. Sans l'hypothèse, deux positions qui se renvoient l'une à l'autre,
l'une dans $\widetilde{G}$ et l'autre dans $\widetilde{P}$, vérifient 1) à 4)
sans que personne ne gagne. La seconde vient de ce que $G \cap P = \emptyset$
quand les $G_0$ sont disjoints.

\textbf{Corollaire} (page 33). Pour un jeu alterné sans partie infinie où
$C_0 = G_0(j) \sqcup G_0(j')$, le couple $(G, P)$ est l'\emph{unique} couple
de parties complémentaires de $C$ qui vérifie 1) à 4).

Dans la convention normale, où les positions terminales sont perdantes, les
conditions disent que $P$ ne contient aucun coup entre deux de ses éléments,
et que toute position hors de $P$ a un coup vers $P$ : $P$ est le
\emph{noyau} du graphe du jeu, et le corollaire est l'existence et l'unicité
du noyau d'un graphe sans chemin infini.\footnote{Von Neumann et Morgenstern
(1944), sous le nom de « solution » d'une relation de domination ; le mot
« noyau » est de Berge (1958). Les P-positions sont aussi les zéros de la
fonction de Sprague et Grundy (1935, 1939). Rien sur la page ne cite ces
travaux.}

\pagerange{47}{47}
\section*{III. Jeux à $N$ joueurs, troisième mouture (page 47)}

La page 47 repart d'un autre cadre. Un jeu à $N$ joueurs est la donnée
\begin{itemize}
  \item d'un ensemble $C$ de configurations et d'un ensemble $J$ de joueurs,
  de cardinal $N$ ;
  \item d'une application $j : C \to J$, définie cette fois partout : en
  toute configuration, il y a un joueur « dont c'est le tour » ;
  \item d'une application $\sigma : C \to \mathfrak{P}(C)$ qui donne les
  configurations suivantes, ou de la relation équivalente
  $\Sigma \subset C \times C$ ;
  \item d'une configuration initiale $x_0$ ; les configurations terminales
  sont celles où $\sigma(x) = \emptyset$.
\end{itemize}
Une partie à $n$ coups est une suite $(x_1, \ldots, x_n)$ avec
$x_i \in \sigma(x_{i-1})$ : le joueur $j(x_0)$ choisit $x_1$, puis $j(x_1)$
choisit $x_2$, et ainsi de suite.\footnote{La lettre $\Sigma$ désigne ici la
relation de succession, que les deux autres moutures notent $R$, et non une
stratégie. La lettre $j$ de l'application est surchargée sur la page.}

Le jeu est \emph{circulaire} si $j$ est surjective (tout joueur joue), s'il
existe $\rho : J \to J$ avec $j(y) = \rho(j(x))$ pour tout coup
$(x, y) \in \Sigma$, et si $\rho$ est une permutation circulaire : les
joueurs jouent à tour de rôle, dans un ordre fixé. L'application $\rho$ est
unique dès que chaque joueur a la main dans au moins une configuration non
terminale.

Trois choses changent par rapport aux deux premières moutures : une position
initiale est fixée, l'initiative est définie aussi aux positions terminales,
comme le $\hat\alpha$ du jeu alterné, et il n'y a pas encore de condition de
gain. Pour $N = 2$, un jeu circulaire est un jeu alterné au sens de la page
31. La page s'arrête aux définitions.

\pagerange{14}{21}
\section*{IV. Réarrangements sur un groupe (pages 14 à 21)}

\pagerange{14}{16}
\subsection*{Les brouillons (pages 14 à 16)}

Les pages 14 et 15 sont des calculs sans phrases, sur des parties finies d'un
groupe commutatif noté additivement : sommes de parties, tableaux de sommes
$0, a, a'$ contre $0, a, b'$, et des listes d'entiers comme $3\,2\,2\,1\,1$ ou
$3\,2\,2\,2$. On y reconnaît --- la page ne le dit pas --- des \emph{profils}
de convolution : la suite décroissante des valeurs de
$\mathbf{1}_E * \mathbf{1}_F$, c'est-à-dire du nombre de façons d'écrire
chaque élément comme somme d'un élément de $E$ et d'un élément de
$F$.\footnote{Lecture nôtre, vérifiée sur deux cas. Pour $E = F = \{0, 1, 2\}$
dans $\mathbb{Z}$, le profil est $(3, 2, 2, 1, 1)$. Pour $E = V_1 \cup \{x\}$,
où $V_1$ est un sous-groupe d'ordre $2$ d'un groupe $G$ d'ordre $4$,
$G = V_1 \cup (x + V_1)$ et $2x = u$, la page 15 calcule
$\mathbf{1}_E * \mathbf{1}_E = 2 \cdot \mathbf{1}_G + \delta_u$, de profil
$(3, 2, 2, 2)$, le nombre qu'elle encadre.} La page 14 s'essaie aussi à
décrire des familles emboîtées $E_2 \subset E_{2n}$ de parties, où $E_2$ est
un sous-groupe et $E_{2n}$ est stable par translation par $E_2$ : ce sont les
candidats au rôle d'« intervalles » dans un groupe qui a de la $2$-torsion.

La page 16, à l'encre bleue, est plus organisée, et ses deux moitiés sont
écrites tête-bêche. L'une est un modèle : dans $G = \mathbb{Z}/(2m+1)$, les
intervalles $E_i = [-i, i]$, de cardinal impair, et
$F_j = [m - j, m + 1 + j]$, de cardinal pair, centrés en $m + \tfrac12$, le
« demi-point » opposé à $0$ sur le cercle (dans $\mathbb{Z}/(2m+1)$, où $2$
est inversible, on a encore $F_j = -F_j$) ; la convolution de deux $E$ ou de deux $F$ est une
combinaison à coefficients positifs des $E_k$, celle d'un $E$ et d'un $F$ une
combinaison des $F_k$.\footnote{La page écrit $E_i = [-i, i]$, le signe moins
tracé à travers le premier $i$, et « NB $E_1 = \varepsilon$ » ; le sens demande
$E_0 = \{0\}$. La lettre des coefficients de la troisième relation est lue
$\gamma$ sans certitude.}
L'autre moitié en tire des axiomes, pour un groupe commutatif $G$ muni d'un
homomorphisme surjectif $\eta : \widetilde{G} \to G$ dont le noyau est
identifié par $i$ au sous-groupe ${}_2G$ des éléments d'ordre divisant $2$.
À chaque $\tilde g \in \widetilde{G}$ est associée une famille
$\mathcal{A}(\tilde g)$ de parties finies de $G$ ; on note
$\mathcal{C}(\tilde g)$ le cône des combinaisons positives de ses éléments,
et $g = \eta(\tilde g)$. Les axiomes sont :
\begin{itemize}
  \item[a)] $\mathcal{A}(\tilde g)$ est totalement ordonnée par inclusion ;
  \item[b)] tout $A \in \mathcal{A}(\tilde g)$ vérifie $A = g - A$ ;
  \item[c)] pour $A \in \mathcal{A}(\tilde g)$ et $B \in \mathcal{A}(\tilde h)$,
  $\mathbf{1}_A * \mathbf{1}_B \in \mathcal{C}(\tilde g \tilde h)$ ;
  \item[d)] pour $u \in {}_2G$, $\mathcal{A}(i(u)\tilde g) = u + \mathcal{A}(\tilde g)$ ;
  \item[d')] pour $\tilde g, \tilde h$, il existe $s \in G$ tel que
  $\mathcal{A}(\tilde g) \cup (s + \mathcal{A}(\tilde h))$ soit totalement
  ordonné.
\end{itemize}
Lu ainsi, $\tilde g$ est le \emph{centre} des ensembles de
$\mathcal{A}(\tilde g)$ : un ensemble tel que $A = g - A$ est symétrique
autour de « $g/2$ », qui n'existe pas toujours dans $G$ --- dans $\mathbb{Z}$,
un intervalle de longueur paire est centré en un demi-entier ---, et
$\widetilde{G}$ est le groupe de ces centres. Deux centres au-dessus du même
$g$ diffèrent d'un élément de ${}_2G$, et d) dit que leurs familles se
déduisent l'une de l'autre par translation.\footnote{Cette lecture est la
nôtre ; la page ne construit pas $\widetilde{G}$, qu'elle se donne. La suite
exacte $0 \to {}_2G \to G \xrightarrow{\,2\,} G$ notée isolément au pied de
la page 21 est la même. La lettre $\eta$ est lue sans certitude, le dernier
crochet reste ouvert, et la page écrit deux conditions d).} C'est le cadre où
le réarrangement « symétrique décroissant » a un sens sur un groupe
commutatif ; il n'est pas repris ailleurs dans le dossier.

\pagerange{18}{18}
\subsection*{Le couplage $|f, g|$ (page 18)}

Soit $G$ un ensemble, $\mathbb{R}[G] = \mathbb{R}^{(G)}$ les fonctions réelles
à support fini, $\mathbb{R}[G]^+$ celles qui sont positives. Deux fonctions
$f, f' \in \mathbb{R}[G]^+$ sont \emph{équivalentes}, $f' \sim f$, si
$f' = f \circ \sigma$ pour une permutation $\sigma$ de $G$ ; il revient au
même de demander qu'elles aient la même suite décroissante de valeurs
\[
f^{\downarrow} = (f^{\downarrow}_0 \geqslant f^{\downarrow}_1 \geqslant \cdots),
\]
valeurs non nulles comptées avec multiplicité et complétées par des zéros, ou
la même image de la mesure de comptage. La page dit $f'$ un
\emph{éparpillement} de $f$.\footnote{La page note $A(f)$ la suite ordonnée,
dans une note de marge lue sans certitude ; on écrit $f^{\downarrow}$, la
notation d'aujourd'hui, qui évite une collision avec les ensembles $A$.} On
pose, pour $f, g \in \mathbb{R}[G]^+$,
\[
|f, g| = \sup_{f' \sim f} \langle f', g \rangle
= \sup_{f' \sim f,\ g' \sim g} \langle f', g' \rangle
= \sum_{n \geqslant 0} f^{\downarrow}_n g^{\downarrow}_n ,
\]
la dernière égalité, que la page n'écrit pas, étant l'inégalité de
réarrangement de Hardy et Littlewood : le couplage est maximal quand on range
les grandes valeurs en face des grandes valeurs. En particulier
$|f, f| = \|f\|_2^2$.

\textbf{Cas d'égalité.} $\langle f', g' \rangle = |f, g|$ si et seulement si,
pour tous $\alpha, \beta > 0$, les ensembles de niveau
$E_\alpha(f') = \{s \mid f'(s) \geqslant \alpha\}$ et $E_\beta(g')$ sont
comparables pour l'inclusion ; autrement dit, si la réunion des familles
$E_*(f')$ et $E_*(g')$ est totalement ordonnée ; autrement dit encore, s'il
existe un ordre total sur $G$ pour lequel $f'$ et $g'$ sont toutes deux
décroissantes.\footnote{La page écrit « $E_\alpha(f') \cup E_\alpha(g')$ soit
tot.\ ordonné » pour la réunion des deux familles, et « $f$ et $g$
décroissantes » pour $f'$ et $g'$. Elle ajoute que l'ordre total peut être
pris de type $[0, n]$ ou $\mathbb{N}$ si $G$ est fini ou dénombrable. On
dit aujourd'hui que $f'$ et $g'$ sont \emph{comonotones}.}

Enfin $f' \sim f$ si et seulement si $|f', g| = |f, g|$ pour tout $g$, et il
suffit de le demander pour les indicatrices des parties finies :
$|f, \mathbf{1}_A|$ est la somme des $\operatorname{card} A$ plus grandes
valeurs de $f$.

\pagerange{18}{19}
\subsection*{La relation $f' \prec f$ (pages 18 et 19)}

\textbf{Proposition.} Pour $f', f \in \mathbb{R}[G]^+$, les conditions
suivantes sont équivalentes :
\begin{itemize}
  \item[a)] $|f', g| \leqslant |f, g|$ pour tout $g \in \mathbb{R}[G]^+$ ;
  \item[b)] $|f', \mathbf{1}_A| \leqslant |f, \mathbf{1}_A|$ pour toute partie
  finie $A$ ;
  \item[c)] $\sum_{i < n} f'^{\downarrow}_i \leqslant \sum_{i < n}
  f^{\downarrow}_i$ pour tout $n \geqslant 1$.
\end{itemize}
On écrit alors $f' \prec f$.\footnote{La page écrit c) avec
$\sum_0^n$ pour $n \in \mathbb{N}^*$, ce qui omet la première inégalité
$f'^{\downarrow}_0 \leqslant f^{\downarrow}_0$ ; on somme sur $i < n$.}

C'est un préordre, dont la relation d'équivalence associée est $\sim$. On
l'appelle aujourd'hui \emph{sous-majorisation faible} : les sommes partielles
de la suite décroissante sont dominées, sans égalité imposée des sommes
totales.\footnote{Hardy, Littlewood et Pólya (1929, puis \emph{Inequalities},
1934) ; l'exposé de référence est Marshall et Olkin, \emph{Inequalities:
Theory of Majorization and Its Applications} (1979). L'équivalence de a) et
c) est une transformation d'Abel, puisque $g^{\downarrow}$ est décroissante et
positive.} Le supremum de $\langle f', g \rangle$ sur les $f' \prec f$ est
$|f, g|$, atteint en des $f' \sim f$ ; et $f' \prec f$ équivaut à
$|f', g'| \leqslant |f, g|$ pour tous $g' \prec g$.

\pagerange{19}{20}
\subsection*{Couples prétassés (pages 19 et 20)}

Désormais $G$ est un groupe, d'élément neutre $e$, et $\mathbb{R}[G]^+$ est
stable par la convolution
$(f * g)(s) = \sum_{ab = s} f(a)\, g(b)$. On note
$\check g(a) = g(a^{-1})$ et $\varepsilon_s$ la masse de Dirac en $s$. On a
\[
(f * g)(s) = \langle f, \varepsilon_s * \check g \rangle ,
\qquad\text{d'où}\qquad
\|f * g\|_\infty \leqslant |f, g| ,
\]
puisque $\varepsilon_s * \check g \sim g$.\footnote{La page écrit
$(f * g)(s) = \langle f * \varepsilon_s, \check g \rangle$ ; ce second membre
vaut $(f * g)(s^{-1})$ dans un groupe commutatif, ce qui est sans effet sur le
supremum en $s$. La forme écrite ici est celle qui rend exacte la condition
d) ci-dessous. La page note l'élément neutre $0$.}

\textbf{Proposition} (page 19). Pour $f, g \in \mathbb{R}[G]^+$, les
conditions suivantes sont équivalentes :
\begin{itemize}
  \item[a)] $\|f' * g'\|_\infty \leqslant \|f * g\|_\infty$ pour tous
  $f' \prec f$, $g' \prec g$ ;
  \item[b)] $(f' * g')(e) \leqslant \|f * g\|_\infty$ pour tous $f' \prec f$,
  $g' \prec g$ (ou seulement $f' \sim f$, $g' \sim g$) ;
  \item[c)] $|f, g| = \|f * g\|_\infty$ ;
  \item[d)] il existe $s \in G$ tel que la réunion des familles $E_*(f)$ et
  $E_*(\varepsilon_s * \check g)$ soit totalement ordonnée.
\end{itemize}
On dit alors que le couple $(f, g)$ est \emph{prétassé}.\footnote{C'est le
mot de la page, dans le titre, « Couples prétassés », dont le premier mot est
surchargé, dans la définition et dans le corollaire. À la page 21, « pré »
est ajouté en interligne au-dessus de « tassé », dans les deux crochets de
la conclusion : il y avait d'abord écrit « tassé ».} La relation n'est pas
symétrique en général, mais l'est si $G$ est commutatif.

\textbf{Corollaire.} Le couple $(f, f)$ est prétassé si et seulement si $f$
est symétrique autour d'un centre : il existe $s \in G$ avec
$\check f = \varepsilon_s * f$, ou, ce qui revient au même,
$\check f = f * \varepsilon_s$.

En effet, deux fonctions équivalentes dont les ensembles de niveau sont
comparables ont les mêmes ensembles de niveau, puisque ceux-ci ont deux à
deux le même cardinal fini : la condition d) force
$f = \varepsilon_s * \check f$.\footnote{L'argument est le nôtre ; la page
énonce le corollaire sans démonstration. L'équivalence des deux formes, à
gauche et à droite, vaut dans tout groupe, avec le même $s$.} Si $g$ est
elle-même symétrique, d) se récrit : il existe $s$ tel que
$E_*(f) \cup E_*(\varepsilon_s * g)$ soit totalement ordonné. Quand $f$ et $g$
sont toutes deux symétriques à translation près, la condition est donc
symétrique en $f$ et $g$ (page 20, en tête), et elle est invariante par
translation.\footnote{Cette remarque de la page 20 est précédée d'un numéro
souligné comme ceux des paragraphes, où se superposent un 4 et un 5 sans
qu'on puisse dire lequel est le dernier écrit ; elle continue pourtant le
paragraphe des couples prétassés, et le §4, « Couples tassés », ne commence
que plus bas sur la page.}

\pagerange{20}{21}
\subsection*{Couples tassés, et la condition K (pages 20 et 21)}

\textbf{Proposition} (page 20). Pour $f, g \in \mathbb{R}[G]^+$, les
conditions suivantes sont équivalentes :
\begin{itemize}
  \item[a)] $f' * g' \prec f * g$ pour tous $f' \prec f$, $g' \prec g$ ;
  \item[b)] $f' * g' \prec f * g$ pour tous $f' \sim f$, $g' \sim g$ ;
  \item[c)] pour tout $h \in \mathbb{R}[G]^+$,
  \[
  \sup_{f' \prec f,\ g' \prec g,\ h' \prec h} (f' * g' * h')(e)
  \leqslant \sup_{h' \sim h} (f * g * h')(e) ;
  \]
  \item[c')] la même inégalité pour $h = \mathbf{1}_A$, $A$ partie finie de
  $G$, et $f' \sim f$, $g' \sim g$, $h' \sim h$.
\end{itemize}
On dit que $(f, g)$ est \emph{tassé}.\footnote{C'est le mot de la page, dans
le titre, « Couples tassés », et au a) du corollaire ; la page 21 le reprend
pour la condition a) sur les parties finies. Le second membre
de c) est aussi $\sup_{h' \sim h} \|f * g * h'\|_\infty$, comme la page
l'écrit. L'équivalence de a) et b) repose sur la description des
sous-majorations faibles par les matrices doublement sous-stochastiques et
sur la sous-additivité des sommes partielles ; la page ne la justifie pas.}

Un couple tassé est prétassé, comme le dit la page : la première inégalité
de $f' * g' \prec f * g$ est
$\|f' * g'\|_\infty \leqslant \|f * g\|_\infty$. La réciproque est
fausse.\footnote{La page écrit, d'une écriture abrégée, « Cela implique que
$(f, g)$ est prétassé ». Contre-exemple (nôtre) à la réciproque :
dans $\mathbb{Z}$, $A = \{0, 1, 3, 4\}$ vérifie $A = 4 - A$, donc
$(\mathbf{1}_A, \mathbf{1}_A)$ est prétassé. Mais $I = \{0, 1, 2, 3\}$ est
équivalent à $A$, et les deux plus grandes valeurs de
$\mathbf{1}_I * \mathbf{1}_I$ sont $4$ et $3$, de somme $7$, contre $4$ et
$2$, de somme $6$, pour $\mathbf{1}_A * \mathbf{1}_A$ ; donc
$\mathbf{1}_I * \mathbf{1}_I \not\prec \mathbf{1}_A * \mathbf{1}_A$.}

\textbf{Corollaire} (page 20). Soit $S \subset \mathbb{R}[G]^+$ stable par
convolution. Les conditions suivantes sont équivalentes :
\begin{itemize}
  \item[a)] tout couple de $S \times S$ est tassé ;
  \item[b)] pour tout $n \geqslant 2$, tous $f_1, \ldots, f_n \in S$ et tous
  $f'_i \prec f_i$,
  $f'_1 * \cdots * f'_n \prec f_1 * \cdots * f_n$ ;
  \item[c)] (\emph{condition K}) si $S$ contient, pour chaque
  $n \leqslant \operatorname{card} G$, l'indicatrice d'une partie finie de
  cardinal $n$ : pour tout $n \geqslant 2$, tous $f_i \in S$ et tous
  $f'_i \prec f_i$,
  \[
  \|f'_1 * \cdots * f'_n\|_\infty \leqslant \|f_1 * \cdots * f_n\|_\infty .
  \]
\end{itemize}
Le passage de a) à b) se fait par récurrence, $S$ étant stable ; celui de c)
à a) passe par c') de la proposition, et c'est là que sert l'hypothèse sur
les cardinaux.\footnote{Cette hypothèse et la borne
$n \leqslant \operatorname{card} G$ sont ajoutées à l'encre bleue, avec le
nom « condition K ». Une condition d), qui ne portait que sur trois
fonctions, est biffée.}

La condition K est la forme que prend sur un groupe l'\emph{inégalité de
Riesz--Sobolev} et sa version à plusieurs fonctions : si $S$ est la classe
des fonctions « symétriques décroissantes », elle dit que la convolution de
fonctions réarrangées n'a jamais un pic plus haut que celle de leurs
réarrangements symétriques décroissants.\footnote{F. Riesz (1930) sur
$\mathbb{R}$, Sobolev (1938) sur $\mathbb{R}^n$, Brascamp, Lieb et Luttinger
(1974) pour plusieurs fonctions ; sur $\mathbb{Z}$, le cas de trois suites
est un théorème de Hardy et Littlewood (\emph{Inequalities}, chap. X), où
deux réarrangements symétriques, l'un centré sur un entier et l'autre sur un
demi-entier, rappellent le revêtement $\widetilde{G}$ de la page 16. Rien sur la page ne cite ces travaux.}

La page 21 reprend la question pour les parties finies. Soit $S$ stable par
convolution et par translation, et $S_0 = S \cap \mathfrak{P}_f(G)$ ses
parties finies (identifiées à leurs indicatrices), supposé non vide. Elle
énonce
\begin{itemize}
  \item[d)] pour $A, B \in S_0$, $A' \sim A$, $B' \sim B$ :
  $\mathbf{1}_{A'} * \mathbf{1}_{B'} \prec \mathbf{1}_A * \mathbf{1}_B$,
  c'est-à-dire que $(A, B)$ est tassé ;
  \item[e)] pour $A, B, C \in S_0$ et des parties équivalentes :
  $\|\mathbf{1}_{A'} * \mathbf{1}_{B'} * \mathbf{1}_{C'}\|_\infty \leqslant
  \|\mathbf{1}_A * \mathbf{1}_B * \mathbf{1}_C\|_\infty$,
\end{itemize}
et conclut que tout $f \in S$ vérifie a) $\check f = \varepsilon_t * f$
pour un $t \in G$ --- $(f, f)$ est prétassé --- et b) ses ensembles de niveau
sont dans $S_0$.\footnote{La page écrit a) « $f = \varepsilon_t * f$ »,
l'accent n'étant pas visible ; le crochet qui suit, « $(f, f)$ prétassé »,
et le corollaire de la page 19 demandent $\check f$. Le premier mot de b) est
illégible, et le crochet de b) dit « $(f, A)$ prétassé pour $A \in S_0$ » ;
que les ensembles de niveau de $f$ soient dans $S_0$ ne découle pas de ce qui
précède, et fonctionne ici comme une définition. La note marginale en tête de
page n'est lue qu'en partie.} Elle définit alors la classe $\widetilde{S}$
des $f$ qui vérifient a) et b), c'est-à-dire
\[
f = \sum_i \alpha_i \mathbf{1}_{A_i},
\qquad \alpha_i > 0,\ \ A_i \in S_0 \ \text{emboîtées},\quad
A_i = s - A_i \ \text{pour un même } s,
\]
calcule $f * g$ pour deux telles fonctions, et veut montrer que, si les
ensembles de niveau des $\mathbf{1}_{A_i} * \mathbf{1}_{B_j}$ forment une
chaîne, ils ont un point commun. La page s'arrête sur cette intersection.
C'est le programme de la page 16 : bâtir, à partir d'une chaîne de parties
symétriques, la classe des fonctions symétriques décroissantes d'un groupe,
et montrer qu'elle est stable par convolution.

Un cas particulier montre ce que la condition d) contient. Pour
$G = \mathbb{Z}/p\mathbb{Z}$, $p$ premier, et deux intervalles de longueurs
$a, b$ avec $a + b \leqslant p + 1$, la condition d) équivaut --- calcul
nôtre --- au théorème de Pollard (1974), qui généralise celui de
Cauchy--Davenport : pour $t \leqslant \min(a, b)$,
$\sum_x \min(r(x), t) \geqslant t(a + b - t)$, où $r(x)$ est le nombre de
représentations de $x$ comme somme d'un élément de $A'$ et d'un élément de
$B'$.\footnote{Les deux fonctions comparées ont la même masse $ab$ ; la
majorisation $r \prec r_0$ équivaut alors à
$\sum_x (r(x) - t)^+ \leqslant \sum_x (r_0(x) - t)^+$ pour tout entier $t$,
soit $\sum_x \min(r(x), t) \geqslant \sum_x \min(r_0(x), t)$, et pour deux
intervalles le second membre vaut $t(a + b - t)$. Pollard, J. London Math.
Soc. (2) 8 (1974) ; Cauchy (1813), Davenport (1935). Rien sur la page ne
cite ces travaux.}

\pagerange{35}{45}
\section*{V. Exponentielles en caractéristique zéro (pages 35 à 45)}

\pagerange{35}{35}
\subsection*{L'exponentielle des éléments nilpotents (page 35)}

Soit $\mathbb{A}$ une $\mathbb{Q}$-algèbre associative unitaire et
$N(\mathbb{A})$ l'ensemble de ses éléments nilpotents ; on a
$1 + N(\mathbb{A}) \subset \mathbb{A}^*$. Sur $N(\mathbb{A})$, la série
$\exp x = \sum_{n \geqslant 0} x^n / n!$ est une somme finie.

\textbf{Théorème 1.}
\begin{itemize}
  \item[1)] $\exp : N(\mathbb{A}) \to 1 + N(\mathbb{A})$ est bijective,
  d'inverse $\log(1 + y) = \sum_{n \geqslant 1} (-1)^{n+1} y^n / n$.
  \item[2)] Si $x, y \in N(\mathbb{A})$ commutent, $x + y \in N(\mathbb{A})$
  et $\exp(x + y) = \exp x \exp y$ ; si $u, v \in 1 + N(\mathbb{A})$
  commutent, $uv \in 1 + N(\mathbb{A})$ et $\log(uv) = \log u + \log v$. On a
  $\exp 0 = 1$ et $\log 1 = 0$.
  \item[3)] Si $L \subset N(\mathbb{A})$ est un sous-$\mathbb{Q}$-espace
  stable par les puissances $x \mapsto x^n$ ($n \geqslant 1$), $\exp$ induit
  une bijection $L \to 1 + L$.
\end{itemize}
Tout se ramène à des identités entre séries formelles à une ou deux
variables qui commutent.\footnote{La page écrit « NB $\exp 1 = 0$ » pour
$\log 1 = 0$.}

\textbf{Remarques.} Si de plus $L$ est stable par produit --- une
sous-algèbre non unitaire formée d'éléments nilpotents ---, $1 + L$ est un
sous-groupe de $\mathbb{A}^*$, et $\exp$ transporte sa loi sur $L$ :
$x \star y = \log(\exp x \exp y)$. Cette loi est commutative si et seulement
si $L$ l'est, et c'est alors l'addition. La page marque la réciproque d'un
« ? » ; elle est vraie.\footnote{Argument nôtre : si
$\exp(\lambda x)\exp(\mu y) = \exp(\mu y)\exp(\lambda x)$ pour tous
$\lambda, \mu \in \mathbb{Q}$, les deux membres sont des polynômes en
$\lambda, \mu$ à coefficients dans $\mathbb{A}$, égaux sur
$\mathbb{Q}^2$, donc coefficient par coefficient ; le coefficient de
$\lambda\mu$ donne $xy = yx$. Quand $L$ est nilpotente, la loi $\star$ est
donnée par la formule de Baker--Campbell--Hausdorff, qui est alors finie ;
c'est le point de départ de la correspondance de Malcev et Lazard (1954)
entre groupes nilpotents uniquement divisibles et algèbres de Lie nilpotentes
sur $\mathbb{Q}$.}

\textbf{Exemple 1.} Soient $K$ une $\mathbb{Q}$-algèbre commutative,
$\mathbb{A}$ une $K$-algèbre et $J$ un nilidéal de $K$. Alors $J\mathbb{A}$
vérifie les conditions de la remarque --- un élément
$\sum j_k a_k$ fait intervenir un nombre fini d'éléments nilpotents et
centraux $j_k$, et est donc nilpotent --- et
$\exp : J\mathbb{A} \to 1 + J\mathbb{A}$ est une bijection. Avec
$K_0 = K/J$ et $\mathbb{A}_0 = \mathbb{A} \otimes_K K_0$, $J\mathbb{A}$ est
le noyau de $\mathbb{A} \to \mathbb{A}_0$ et $1 + J\mathbb{A}$ l'image
réciproque de $1$.

\textbf{Exemple 2.} Soit $M$ un $K$-module, $M_0 = M/JM$, et
$\mathcal{U} = \operatorname{End}_K(M)$. On pose
\[
L = \{u \in \mathcal{U} \mid u(M) \subset JM\},
\qquad
1 + L = \{u \in \mathcal{U} \mid u \otimes_K K_0 = \mathrm{id}_{M_0}\}.
\]
Si $J$ est \emph{nilpotent}, ou si $M$ est de type fini, $L$ vérifie les
conditions de la remarque et $\exp : L \to 1 + L$ est une
bijection.\footnote{L'hypothèse n'est pas sur la page, qui dit que « les
conditions de la remarque sont satisfaites » pour un nilidéal. Une ligne
illisible de l'exemple 1 commence « mais avec $J$ nilpotent ». Si $J$ est
seulement nil et $M$ quelconque, un $u \in L$ peut ne pas être nilpotent
(exemple nôtre) : $K = \mathbb{Q}[t_1, t_2, \ldots]/(t_1^2, t_2^2, \ldots)$,
$J = (t_1, t_2, \ldots)$, $M = K^{(\mathbb{N})}$ de base $(e_k)$, et
$u(e_k) = t_k e_{k+1}$ ; alors $u^n(e_1) = t_1 \cdots t_n e_{n+1} \neq 0$.
Pour $M$ de type fini, les coefficients de $u$ sur des générateurs engendrent
un idéal de type fini formé de nilpotents, donc nilpotent.}

\pagerange{36}{37}
\subsection*{Groupes formels à un paramètre (pages 36 et 37)}

Soit $K$ une $\mathbb{Q}$-algèbre commutative unitaire. Les foncteurs
considérés sont définis sur les $K$-algèbres commutatives $K'$. La
\emph{droite formelle} $\widehat{\mathbb{G}}_a$ est le sous-foncteur
$K' \mapsto N(K')$ du groupe additif $\mathbb{G}_a : K' \mapsto K'$ ; c'est
la limite inductive des $\operatorname{Spec} K[t]/(t^n)$.\footnote{La page
note $\overline{\mathbb{G}}_{a,K}$ la droite formelle ; on écrit
$\widehat{\mathbb{G}}_a$, la notation d'aujourd'hui du groupe additif
formel.} Pour une $K$-algèbre $\mathcal{U}$, soit $W(\mathcal{U})$ le
foncteur $K' \mapsto \mathcal{U} \otimes_K K'$ et
$\mathrm{Un}(\mathcal{U}) \subset W(\mathcal{U})$ le foncteur en groupes
$K' \mapsto (\mathcal{U} \otimes_K K')^*$.\footnote{La lettre $W$ est lue
sans certitude à la page 37. Quand $\mathcal{U}$ est un $K$-module projectif
de type fini, $\mathrm{Un}(\mathcal{U})$ est représentable, et c'est le
schéma en groupes des unités de $\mathcal{U}$ ; la page le traite en
foncteur.}

À $x \in \mathcal{U}$ on associe $\bar e_x : \widehat{\mathbb{G}}_a \to
\mathrm{Un}(\mathcal{U})$, $\lambda \mapsto \exp(\lambda x)$, bien défini
parce que $\lambda x$ est nilpotent quand $\lambda$ l'est (exemple 1), et
homomorphisme par le 2) du théorème 1. Si $x$ est nilpotent,
$e_x : \mathbb{G}_a \to \mathrm{Un}(\mathcal{U})$,
$\lambda \mapsto \exp(\lambda x)$, est défini pour tout $\lambda$, et
$\bar e_x$ en est la restriction. D'où deux applications
\[
(*)\quad \mathcal{U} \to
\operatorname{Hom}_{\mathrm{gr}}(\widehat{\mathbb{G}}_a, \mathrm{Un}(\mathcal{U})),
\ \ x \mapsto \bar e_x ;
\qquad
(**)\quad N(\mathcal{U}) \to
\operatorname{Hom}_{\mathrm{gr}}(\mathbb{G}_a, \mathrm{Un}(\mathcal{U})),
\ \ x \mapsto e_x .
\]
Si $x$ et $y$ commutent, $\bar e_{x+y} = \bar e_x \bar e_y$, et de même pour
$e$ ; si $\mathcal{U}$ est commutative, $(*)$ et $(**)$ sont des
homomorphismes de groupes.

\textbf{Théorème 2.} $(*)$ et $(**)$ sont bijectives.

Un morphisme de foncteurs $\widehat{\mathbb{G}}_a \to W(\mathcal{U})$ est une
série $\sum_{n \geqslant 0} a_n \lambda^n$ à coefficients $a_n \in
\mathcal{U}$, et un morphisme $\mathbb{G}_a \to W(\mathcal{U})$ un polynôme :
c'est le lemme de Yoneda, $W(\mathcal{U})$ prenant la valeur
$\mathcal{U}[t]/(t^n)$ sur $K[t]/(t^n)$ et $\mathcal{U}[t]$ sur
$K[t]$.\footnote{La page utilise cette description sans la justifier.}
L'injectivité se lit sur le coefficient de $\lambda$. Pour la surjectivité,
écrire que $\varphi(\lambda + \mu) = \varphi(\lambda)\varphi(\mu)$ donne,
coefficient par coefficient,
\[
\binom{i + j}{i}\, a_{i+j} = a_i a_j ,
\qquad\text{d'où}\qquad
n a_n = a_1 a_{n-1}
\quad\text{et}\quad
a_n = \frac{x^n}{n!},\ \ x = a_1 ,
\]
où l'on divise par $n$ : c'est là que sert la caractéristique zéro. Pour
$\mathbb{G}_a$, les $a_n$ sont nuls à partir d'un certain rang, donc $x$ est
nilpotent.\footnote{Les homomorphismes de $\mathbb{G}_a$ vers le groupe des
unités d'une algèbre sont donc les $\lambda \mapsto \exp(\lambda x)$, $x$
nilpotent, et ceux de $\widehat{\mathbb{G}}_a$ les mêmes avec $x$
quelconque : en caractéristique zéro, un groupe formel est déterminé par son
algèbre de Lie (Cartier), et l'algèbre de Lie de $\mathrm{Un}(\mathcal{U})$
est $\mathcal{U}$. Voir Demazure et Gabriel, \emph{Groupes algébriques}
(1970), et SGA 3, que Grothendieck a codirigé ; rien sur la page ne les
cite.}

\pagerange{38}{38}
\subsection*{Opérations sur un module (page 38)}

Soit $M$ un $K$-module et $\mathbf{Aut}_K(M)$ le foncteur en groupes
$K' \mapsto \operatorname{Aut}_{K'}(M \otimes_K K')$. Le même calcul
donne :
\begin{itemize}
  \item[$(*)$] $u \mapsto \bar e_u$ est une bijection de
  $\operatorname{End}_K(M)$ sur
  $\operatorname{Hom}_{\mathrm{gr}}(\widehat{\mathbb{G}}_a, \mathbf{Aut}_K(M))$,
  pour tout $M$ ;
  \item[$(**)$] $u \mapsto e_u$ est une bijection de l'ensemble des
  endomorphismes \emph{localement nilpotents} de $M$ (tout $m \in M$ est tué
  par une puissance de $u$) sur
  $\operatorname{Hom}_{\mathrm{gr}}(\mathbb{G}_a, \mathbf{Aut}_K(M))$ ; si $M$
  est de type fini, ce sont les endomorphismes nilpotents.
\end{itemize}
Un morphisme de foncteurs $\widehat{\mathbb{G}}_a \to \mathbf{End}_K(M)$
s'écrit bien de façon unique $\sum_n u_n \lambda^n$ avec
$u_n \in \operatorname{End}_K(M)$, puisque
$\operatorname{End}_{K[t]/(t^n)}(M[t]/(t^n)) = \operatorname{End}_K(M)^n$.
Pour $\mathbb{G}_a$, en revanche, un morphisme est une application
$K$-linéaire $M \to M[t]$, c'est-à-dire une suite $(u_n)$ telle que, pour
chaque $m$, les $u_n(m)$ soient nuls pour $n$ grand --- non une suite
d'endomorphismes nuls à partir d'un certain rang, comme l'écrit la page ; les
deux coïncident si $M$ est de type fini.\footnote{Contre-exemple (nôtre) à
l'énoncé de la page : $M = K^{(\mathbb{N})}$ de base $(e_k)$, $u(e_0) = 0$,
$u(e_k) = e_{k-1}$. L'endomorphisme $u$ est localement nilpotent sans être
nilpotent, et $\lambda \mapsto \exp(\lambda u)$ est un homomorphisme
$\mathbb{G}_a \to \mathbf{Aut}_K(M)$ qui n'est pas de la forme $e_x$ avec
$x$ nilpotent. Qu'en caractéristique zéro les actions de $\mathbb{G}_a$
correspondent aux endomorphismes, et pour une algèbre aux dérivations,
localement nilpotents est un résultat classique. La page note elle-même que
le cas projectif de type fini se ramène au précédent avec
$\mathcal{U} = \operatorname{End}_K(M)$.}

\pagerange{39}{41}
\subsection*{Morphismes, et le passage à $\exp$ (pages 39 à 41)}

\textbf{Proposition} (page 39). Soient $u \in \operatorname{End}_K(M)$,
$v \in \operatorname{End}_K(N)$ et $f : M \to N$ linéaire. Pour que $f$
commute aux actions $\bar e_u$ et $\bar e_v$ de $\widehat{\mathbb{G}}_a$, il
faut et il suffit que $vf = fu$ ; si $u$ et $v$ sont nilpotents, c'est aussi
la condition pour $\mathbb{G}_a$.

On écrit $\sum_n \lambda^n (v^n f - f u^n)/n! = 0$ identiquement en
$\lambda$ nilpotent, ce qui équivaut à $v^n f = f u^n$ pour tout $n$, et donc
à $vf = fu$.

\textbf{Proposition} (pages 39 à 41). Soient $u, v$ nilpotents. Les
conditions suivantes sont équivalentes : (1) $f$ est un
$\widehat{\mathbb{G}}_a$-morphisme ; (2) $f$ est un
$\mathbb{G}_a$-morphisme ; (3) $\delta = vf - fu$ est nul ; (4)
$f \circ \exp u = \exp v \circ f$.

Seul (4) $\Rightarrow$ (3) est nouveau. Avec
$\delta_n = v^n f - f u^n = \sum_{i + j = n - 1} v^i \delta u^j$, la
condition (4) s'écrit $\sum_{n \geqslant 1} \delta_n / n! = 0$, soit
\[
\delta = \sum_{i + j \geqslant 1} c_{ij}\, v^i \delta u^j,
\qquad c_{ij} = -\frac{1}{(i + j + 1)!} .
\]
Les opérateurs $\Theta(x) = vx$ et $\Theta'(x) = xu$ sur
$\operatorname{Hom}_K(M, N)$ sont nilpotents et commutent ; donc
$\rho = \sum_{i+j \geqslant 1} c_{ij} \Theta^i \Theta'^j$ est nilpotent, et
$\delta = \rho\delta = \rho^N \delta = 0$ pour $N$ grand (page
41).\footnote{La page écrit $c_{ij} = 1/(i+j)!$ ; l'équation qui précède
donne la valeur ci-dessus, et seule compte l'existence des $c_{ij}$. Elle
fait agir $\Theta$ et $\Theta'$ sur $\mathcal{U} = \operatorname{End}_K(M)$ ;
$\delta$ vit dans $\operatorname{Hom}_K(M, N)$, et c'est là qu'ils opèrent.
Une voie plus courte (nôtre) : de $f a = b f$ on tire $f P(a) = P(b) f$ pour
tout polynôme $P$, et $\log$ est un polynôme sur les unipotents d'ordre
donné ; appliqué à $a = \exp u$, $b = \exp v$, cela donne $fu = vf$.}

\pagerange{41}{42}
\subsection*{Tenseurs, dérivations, et l'action sur $\operatorname{Hom}$ (pages 41 et 42)}

Si $\widehat{\mathbb{G}}_a$ opère sur $M$ par $\bar e_u$ et sur $N$ par
$\bar e_v$, il opère sur $M \otimes N$ par
\[
\bar e_u \otimes \bar e_v = \bar e_{u \otimes 1 + 1 \otimes v},
\]
puisque $u \otimes 1$ et $1 \otimes v$ commutent ; de même pour $e$ si $u$ et
$v$ sont nilpotents, et pour un produit tensoriel fini quelconque.

\textbf{Corollaire.} Soit $M$ une $K$-algèbre, non nécessairement associative
ni unitaire, et $u \in \operatorname{End}_K(M)$. L'action $\bar e_u$ de
$\widehat{\mathbb{G}}_a$ se fait par automorphismes d'algèbre si et seulement
si $u$ est une dérivation. Si $u$ est nilpotent, il en est de même de
l'action $e_u$ de $\mathbb{G}_a$, et aussi de la seule condition que
$\exp u$ soit un automorphisme d'algèbre.

On applique la proposition précédente à la multiplication
$m : M \otimes M \to M$ : elle est équivariante si et seulement si
$u \circ m = m \circ (u \otimes 1 + 1 \otimes u)$, ce qui est la règle de
Leibniz ; et la condition sur $\exp u$ est la condition (4).\footnote{La page
qualifie ce corollaire de « naïf », mot d'une lecture incertaine.}

Si $\mathbb{G}_a$ opère sur $M$ et $N$ par $e_u$ et $e_v$, il opère sur
$\operatorname{Hom}_K(M, N)$ par $e_\delta$, avec $\delta(f) = vf - fu$, et de
même pour $\widehat{\mathbb{G}}_a$. La page note qu'en toute rigueur ces
actions ne sont définies que si le foncteur
$K' \mapsto \operatorname{Hom}_{K'}(M \otimes K', N \otimes K')$ est de la
forme $W(\operatorname{Hom}_K(M, N))$, ce qui n'est guère assuré que si $M$
est projectif de type fini.

\pagerange{42}{43}
\subsection*{Foncteurs en modules (pages 42 et 43)}

Pour s'affranchir de cette restriction, la page généralise aux
\emph{$K$-foncteurs en modules} $\mathcal{M}$ : pour chaque $K'$, un
$K'$-module $\mathcal{M}(K')$, et pour chaque $K' \to K''$ une application
semi-linéaire $\mathcal{M}(K') \to \mathcal{M}(K'')$.\footnote{La lettre qui
qualifie ces modules est surchargée et entourée ; la transcription lit
$\mathbb{Q}$, et le contexte attend $K$.} Les endomorphismes de
$\mathcal{M}$ forment une $K$-algèbre. Si $u$ en est un élément nilpotent,
les $\exp(\lambda u)$, $\lambda \in K'$, font de $\mathcal{M}$ un
$\mathbb{G}_a$-module ; pour $u$ quelconque,
$\bar e_u(\lambda) = \exp(\lambda u)$, $\lambda$ nilpotent, définit un
homomorphisme $\widehat{\mathbb{G}}_a \to \mathbf{Aut}_K(\mathcal{M})$. La
page se demande, sans y répondre, si
\[
\operatorname{End}_K(\mathcal{M}) \to
\operatorname{Hom}_{\mathrm{gr}}(\widehat{\mathbb{G}}_a, \mathbf{Aut}_K(\mathcal{M})),
\qquad
N(\operatorname{End}_K(\mathcal{M})) \to
\operatorname{Hom}_{\mathrm{gr}}(\mathbb{G}_a, \mathbf{Aut}_K(\mathcal{M}))
\]
sont encore bijectives, et annonce l'extension aux algèbres symétriques,
tensorielles et extérieures. On ne tranche pas la question ; la page 38 montre
seulement que, pour $\mathcal{M} = W(M)$, la seconde demande de remplacer les
nilpotents par les localement nilpotents.

\pagerange{43}{43}
\subsection*{Globalisation (page 43)}

Soit $X$ un schéma de caractéristique zéro et $\Theta$ une dérivation de
$\mathcal{O}_X$, nilpotente localement sur $X$ : tout point a un voisinage
ouvert sur lequel une puissance de $\Theta$ est nulle. Alors $\exp \Theta$ est
un automorphisme du faisceau d'anneaux $\mathcal{O}_X$ qui induit l'identité
sur $\mathcal{O}_{X_{\mathrm{red}}}$, donc un automorphisme de $X$, qu'on note
$\mathbb{E}_\Theta$, et qui est un $S$-automorphisme si $\Theta$ est une
$S$-dérivation. On a $\mathbb{E}_0 = \mathrm{id}$ et
\[
\mathbb{E}_{\Theta + \Theta'} = \mathbb{E}_\Theta \circ \mathbb{E}_{\Theta'}
\qquad\text{si } [\Theta, \Theta'] = 0 .
\]
Que $\exp \Theta$ soit l'identité modulo les nilpotents tient à ce qu'une
dérivation nilpotente d'une $\mathbb{Q}$-algèbre commutative prend ses
valeurs dans le nilradical.\footnote{Argument nôtre. Une dérivation d'une
$\mathbb{Q}$-algèbre préserve le nilradical, et passe donc à l'anneau réduit.
Sur un anneau réduit, si $\Theta^n = 0$ avec $n$ minimal et $n \geqslant 2$,
soit $b = \Theta^{n-2}(a)$ avec $c = \Theta(b) \neq 0$ ; comme
$\Theta(c) = 0$, on a $0 = \Theta^n(b^n) = n!\, c^n$, d'où $c = 0$. La page
écrit « $\Theta$ loc.\ nilpotente » ; au sens d'aujourd'hui (chaque section
est tuée par une puissance de $\Theta$), l'énoncé serait faux : sur la droite
affine, $d/dx$ est localement nilpotente, et $\exp(\lambda\, d/dx)$ est une
translation, qui n'induit pas l'identité. La page écrit
$\mathbb{E}_\Theta + \mathbb{E}_{\Theta'}$, avec un signe $+$ ; il s'agit de
la composition. L'ordre des facteurs est sans importance ici, puisque les
deux automorphismes commutent.}

\pagerange{44}{45}
\subsection*{Automorphismes unipotents et dérivations nilpotentes (pages 44 et 45)}

Soit $A$ une $\mathbb{Q}$-algèbre commutative unitaire, $u$ un automorphisme
de $A$, et $J$ l'idéal engendré par les $u(x) - x$.

\textbf{Proposition} (page 44). Soit $n \in \mathbb{N}$.
\begin{itemize}
  \item[1)] Les conditions suivantes sont équivalentes : (i) $u$ est un
  opérateur différentiel d'ordre $\leqslant n$ ; (ii) $J^{n+1} = 0$ ; (iii)
  il existe un idéal $I$ avec $I^{n+1} = 0$ et $u \equiv \mathrm{id}$ modulo
  $I$.
  \item[2)] $\Theta \mapsto \exp \Theta$ est une bijection de l'ensemble des
  dérivations $\Theta$ de $A$ telles que $\Theta^{n+1} = 0$ sur l'ensemble des
  automorphismes $u$ tels que $(u - \mathrm{id})^{n+1} = 0$.
\end{itemize}
Pour (i) $\Leftrightarrow$ (ii) : au sens de EGA IV (16.8), un opérateur $P$
est différentiel d'ordre $\leqslant n$ si ses commutateurs itérés
$[\ldots[[P, a_0], a_1], \ldots, a_n]$ avec les multiplications sont nuls ;
pour un endomorphisme d'anneau, $[u, a] = (u(a) - a)\, u$, et le commutateur
itéré vaut $\prod_{i=0}^n (u(a_i) - a_i) \cdot u$, qui est nul pour tous les
$a_i$ si et seulement si $J^{n+1} = 0$.\footnote{La lecture « op.\ diff. »
est incertaine aux pages 44 et 45 ; ce calcul, qui est le nôtre, la
confirme.} Pour 2), $\log u$ est une dérivation parce que
$\log(u \otimes u) = \log u \otimes 1 + 1 \otimes \log u$, par l'argument du
corollaire de la page 41.

Une quatrième condition, (iv) $u = \exp \Theta$ avec $\Theta$ dérivation et
$\Theta^{n+1} = 0$, est écrite puis biffée, avec en marge « est-il
nécessaire ? ». La réponse est non. La condition (iv) entraîne (i), parce que
$\Theta^k$ est un opérateur différentiel d'ordre $\leqslant k$ ; mais elle ne
lui est pas équivalente. Dans $A = \mathbb{Q}[x, e]/(e^2)$, l'automorphisme
$u$ qui fixe $x$ et envoie $e$ sur $2e$ vérifie $J = (e)$, $J^2 = 0$, alors
que $(u - \mathrm{id})^k(e) = e$ pour tout $k$ : $u$ n'est pas unipotent, et
n'est l'exponentielle d'aucune dérivation nilpotente.\footnote{Exemple nôtre.
L'interligne au-dessus de (iv) dit « Pour ceci, il suffit que
$(u - \mathrm{id})^{n+1} = 0$ » : c'est la condition de 2), qui est la bonne.
Deux conditions (ii') et (iii') de la marge gauche, portant sur
$\mathrm{gr}_I(A)$, sont biffées et encadrées ; on les retrouve plus bas.}

\textbf{Corollaire} (page 44). Pour un automorphisme $u$ de $A$, les
conditions (i) $u$ est un opérateur différentiel, (ii) $J$ est nilpotent,
(iii) $u \equiv \mathrm{id}$ modulo un idéal nilpotent, sont équivalentes.
La condition (iv) $u = \exp \Theta$ pour une dérivation nilpotente $\Theta$,
c'est-à-dire $u$ \emph{unipotent}, les entraîne et est strictement plus
forte ; et $\Theta \mapsto \exp \Theta$ est une bijection des dérivations
nilpotentes sur les automorphismes unipotents.\footnote{La page énonce
(i) à (iv) comme équivalentes, et la bijection comme allant sur les
automorphismes qui sont des opérateurs différentiels. L'exemple ci-dessus
montre que ces deux énoncés sont faux ; on donne ce qui est vrai.}

\textbf{Corollaire} (page 45). Si le nilradical $N(A)$ est nilpotent (par
exemple si $A$ est noethérien), les conditions (i) à (iii) équivalent à : (v)
$u$ induit l'identité sur $A_{\mathrm{red}}$ ; et si $N(A)^{N+1} = 0$, $u$ est
alors un opérateur différentiel d'ordre $\leqslant N$.

\textbf{Corollaire} (page 45). Si $A$ est une $K$-algèbre, $K$ une
$\mathbb{Q}$-algèbre, et $\Theta$ une dérivation nilpotente, $\exp \Theta$
est $K$-linéaire si et seulement si $\Theta$ l'est.

\textbf{Corollaire} (page 45). Soit $I$ un idéal de $A$ avec $I^{N+1} = 0$.
Soient $L_K(A, I)$ l'ensemble des $K$-dérivations $\Theta$ telles que
$\Theta(A) \subset I$ et $\Theta(I) \subset I^2$, et $G_K(A, I)$ le groupe
des $K$-automorphismes qui induisent l'identité sur
$\mathrm{gr}_I(A) = \bigoplus_k I^k / I^{k+1}$, c'est-à-dire sur $A/I$ et sur
$I/I^2$. Alors $\Theta^i(A) \subset I^i$ pour $\Theta \in L_K(A, I)$, en
particulier $\Theta^{N+1} = 0$, et
\[
\exp : L_K(A, I) \xrightarrow{\ \sim\ } G_K(A, I).
\]
La page définit $L_K(A, I)$ et $G_K(A, I)$ par les seules conditions
$\Theta(A) \subset I$ et $u \equiv \mathrm{id}$ modulo $I$ ; avec ces
définitions, le corollaire est faux.\footnote{Le même anneau
$A = \mathbb{Q}[x, e]/(e^2)$, $I = (e)$, $N = 1$, en donne les
contre-exemples (nôtres) : la dérivation $\Theta$ définie par
$\Theta(x) = e$ et $\Theta(e) = e$ envoie $A$ dans $I$ mais
$\Theta^k(x) = e$ pour tout $k \geqslant 1$, et l'automorphisme $e \mapsto
2e$ induit l'identité sur $A/I$ sans être unipotent. La condition
supplémentaire est celle de la note marginale (iii') de la page 44, biffée :
« $u$ induise l'identité sur $\mathrm{gr}_I(A)$ ». Avec elle, $u -
\mathrm{id}$ envoie $I^k$ dans $I^{k+1}$, donc
$(u - \mathrm{id})^{N+1} = 0$, et la proposition 2) s'applique.}
C'est, dans ce cadre, la correspondance entre une algèbre de Lie nilpotente
de dérivations et un groupe unipotent d'automorphismes : $L_K(A, I)$ est une
algèbre de Lie, filtrée par les $I^k$, et $\exp$ en fait un groupe par la
formule de Baker--Campbell--Hausdorff.

\pagerange{49}{53}
\section*{VI. L'orientation des espaces vectoriels réels (pages 49 à 53)}

\pagerange{49}{49}
\subsection*{Deux lemmes sur les ensembles à deux éléments (page 49)}

Un ensemble $P$ à deux éléments a une seule involution sans point fixe,
$\tau_P$ ; c'est un torseur sous $\mathbb{Z}/2$.\footnote{La page note cette
involution $a_P$, soulignée.}

\textbf{Lemme 1.} Pour $\varphi : P \to Q$ entre ensembles à deux éléments,
sont équivalents : $\varphi$ bijective, injective, surjective, non constante,
$\mathbb{Z}/2$-équivariante ($\varphi \circ \tau_P = \tau_Q \circ \varphi$).

\textbf{Lemme 2.} Pour $\varphi : P \times Q \to R$ entre ensembles à deux
éléments, $\varphi$ se factorise par un isomorphisme
$P \wedge Q \xrightarrow{\sim} R$ si et seulement si, pour $x \in P$,
$\varphi(x, \cdot) : Q \to R$ est bijective et
$\varphi(\tau_P x, \cdot) = \tau_R \circ \varphi(x, \cdot)$. Ici
$P \wedge Q = (P \times Q)/(\mathbb{Z}/2)$ est le produit contracté des deux
torseurs.

\pagerange{49}{49}
\subsection*{Les axiomes (page 49)}

Soit $\mathcal{V}$ le groupoïde des espaces vectoriels réels de dimension
finie et de leurs isomorphismes linéaires. Une \emph{théorie de
l'orientation} est la donnée
\begin{itemize}
  \item[a)] d'un foncteur $\Omega$ de $\mathcal{V}$ vers les ensembles, avec
  $\operatorname{card} \Omega(V) = 2$ pour tout $V$ ;
  \item[b)] d'une identification $\Omega(0) \simeq \{\pm 1\}$ ;
  \item[c)] pour tout hyperplan $W \subset V$ et tout demi-espace fermé $V'$
  de $V$ bordé par $W$, d'une \emph{application de Stokes}
  $\mathrm{St}_{V'} : \Omega(V) \to \Omega(W)$, naturelle en $(V, V')$,
\end{itemize}
soumise à deux axiomes : $\mathrm{St}_{V'}$ est bijective, et si $V''$ est le
demi-espace opposé, $\mathrm{St}_{V''} = \tau \circ \mathrm{St}_{V'}$.\footnote{La page écrit $\mathrm{St}_{V''} = a \circ \mathrm{St}_V$ ; on lit
$\mathrm{St}_{V'}$. Le premier axiome est d'abord justifié par le seul
cardinal, ce qui ne suffit pas : une application entre deux ensembles à deux
éléments peut être constante. Le mot « vectoriels » est d'une lecture
incertaine.} Par le lemme 2, il revient au même de demander que, en notant
$D(W, V)$ l'ensemble à deux éléments des demi-espaces bordés par $W$, les
applications de Stokes se factorisent en un isomorphisme
naturel\footnote{La page note cet ensemble $\mathrm{Dem}_{W,V}$, lecture
incertaine sous une rature, puis $\Omega(W, V)$.}
\[
D(W, V) \wedge \Omega(V) \xrightarrow{\ \sim\ } \Omega(W).
\]
Le nom rappelle la formule de Stokes,
où l'orientation d'un domaine induit celle de son bord.

Un morphisme de théories $(\Omega, \mathrm{St}) \to (\Omega', \mathrm{St}')$
est une transformation naturelle $\Omega \to \Omega'$ compatible à b) et à c).

\textbf{Théorème} (page 49). a) Tout morphisme de théories est un
isomorphisme. b) Entre deux théories de l'orientation, il existe un et un
seul isomorphisme. c) Il existe une théorie de l'orientation.

Pour a) et l'unicité dans b), on raisonne par récurrence sur la dimension :
$\varphi_V$ est déterminé, et bijectif, par
$\varphi_W \circ \mathrm{St}_{V'} = \mathrm{St}'_{V'} \circ \varphi_V$,
puisque les applications de Stokes sont bijectives. L'existence dans b) et c)
passe par une description des théories en termes de groupes.

\pagerange{49}{51}
\subsection*{Réduction au signe du déterminant (pages 49 et 51)}

Le groupoïde des espaces de dimension $n$ est équivalent au groupe
$\mathrm{GL}(n, \mathbb{R})$. Se donner $\Omega$ en dimension $n$ revient
donc à se donner l'ensemble $\omega_n = \Omega(\mathbb{R}^n)$ et l'action de
$\mathrm{GL}(n, \mathbb{R})$ sur lui, c'est-à-dire un homomorphisme
$\varepsilon_n : \mathrm{GL}(n, \mathbb{R}) \to \{\pm 1\}$. Se donner les
applications de Stokes en dimension $n$ revient, en choisissant le
demi-espace $\mathbb{R}^{n-1} \times \mathbb{R}_+$ de $\mathbb{R}^n$, à se
donner une bijection $s_n : \omega_n \to \omega_{n-1}$ compatible à l'action
du stabilisateur de l'hyperplan $\mathbb{R}^{n-1}$, formé des matrices
\[
u_n = \begin{pmatrix} u_{n-1} & * \\ 0 & \lambda \end{pmatrix},
\qquad u_{n-1} \in \mathrm{GL}(n-1, \mathbb{R}),\ \lambda \in \mathbb{R}^* ,
\]
qui échange ou non les deux demi-espaces selon le signe de $\lambda$. La
compatibilité s'écrit
\[
\varepsilon_n(u_n) = \operatorname{sg}(\lambda)\, \varepsilon_{n-1}(u_{n-1}).
\]
Pour $n = 1$, elle donne $\varepsilon_1 = \operatorname{sg}$, et de proche en
proche\footnote{Le passage du stabilisateur à $\mathrm{GL}(n, \mathbb{R})$ tout
entier, que la page fait « de proche en proche », utilise que tout
homomorphisme $\mathrm{GL}(n, \mathbb{R}) \to \{\pm 1\}$ se factorise par le
déterminant, le groupe dérivé étant $\mathrm{SL}(n, \mathbb{R})$ ; il est
donc déterminé par ses valeurs sur les matrices diagonales
$\operatorname{diag}(1, \ldots, 1, \lambda)$. La note marginale, lue
« choisi $\mathbb{R}^{n-1} \times \mathbb{R}^+$ comme demi-espace de
$\mathbb{R}^n$ », a son premier mot incertain. La page 49 écrit la donnée de
Stokes $\Omega(V, W) \wedge \Omega(W) \to \Omega(W)$ ; le second facteur est
$\Omega(V)$.}
\[
\varepsilon_n = \operatorname{sg} \circ \det .
\]

Les caractères sont donc imposés, et une théorie se réduit aux bijections
$s_n : \omega_n \to \omega_{n-1}$ et à $\omega_0 \simeq \{\pm 1\}$, qui
identifient tous les $\omega_n$ à $\{\pm 1\}$ : d'où l'unicité à isomorphisme
unique près, et l'existence en prenant pour $\Omega(V)$ l'ensemble des deux
orientations usuelles de $V$. C'est ce que disait un passage encadré et
biffé d'une grande croix, dont la conclusion est « donc
$\omega_n \simeq \{\pm 1\}$ ».

Dans le langage d'aujourd'hui, $\Omega(V)$ est le \emph{torseur
d'orientation} de $V$, l'ensemble des deux composantes connexes de
$\Lambda^{\dim V} V \smallsetminus \{0\}$, où $\Lambda^{\dim V} V$ est la
\emph{droite déterminant} ; $\operatorname{sg} \circ \det$ est le caractère
d'orientation. Le théorème dit que ces objets sont caractérisés par leur
comportement vis-à-vis des bords. La normalisation de la page, où
$[e_1 \wedge \cdots \wedge e_n]$ s'envoie sur
$[e_1 \wedge \cdots \wedge e_{n-1}]$ pour le demi-espace
$\mathbb{R}^{n-1} \times \mathbb{R}_+$, met en dernier la normale
\emph{rentrante} ; la convention usuelle des géomètres, normale
\emph{sortante} en premier, en diffère par le signe $(-1)^{\dim V}$, et le
théorème fournit l'isomorphisme canonique entre les deux.\footnote{La
normalisation n'est pas écrite à la page 51 ; on la lit dans la formule de
la page 53. Le calcul du signe est le nôtre.}

\pagerange{53}{53}
\subsection*{Produits, et composition des applications de Stokes (page 53)}

La page 53 ébauche la suite, sans phrases. Pour un sous-espace $W \subset V$
quelconque, les applications de Stokes se généralisent en un isomorphisme
$\Omega(V/W) \wedge \Omega(V) \simeq \Omega(W)$ ; pour un drapeau
$U \subset W \subset V$, elles se composent, et la page nomme
$I(\omega, \omega')$ un élément de $\Omega$ qu'elle ne définit pas. Il y a un
produit
\[
\Omega(V) \wedge \Omega(V') \to \Omega(V \times V'),
\qquad
(\omega, \omega') \mapsto \omega \times \omega',
\qquad
\omega' \times \omega = (-1)^{\dim V \cdot \dim V'}\, \omega \times \omega' ,
\]
la seconde égalité s'entendant via l'échange des facteurs : c'est la
\emph{règle des signes de Koszul}. Enfin la page compare Stokes et produits.
Dans la normalisation qu'on vient de dire, pour un demi-espace de $V'$ et pour
un demi-espace de $V$ respectivement,
\[
\mathrm{St}(\omega \times \omega') = \omega \times \mathrm{St}(\omega'),
\qquad
\mathrm{St}(\omega \times \omega') = (-1)^{\dim V'}\, \mathrm{St}(\omega) \times \omega' ,
\]
ce qui est la formule
$\mathrm{St}(e_1 \cdots e_{p-1} e_p f_1 \cdots f_q) =
(-1)^q\, e_1 \cdots e_{p-1} f_1 \cdots f_q$ écrite au bas de la page, où
$e_p$ est la normale.\footnote{La page écrit les deux formules sans signe,
avec au-dessus de la première, relié par un trait, « $(-1)^?$ » ; le signe
est déterminé ici, dans la normalisation de la page 51. Deux croquis au bas de
la page --- un carré, et un quadrant hachuré avec ses demi-axes fléchés et
une flèche de rotation --- illustrent le cas de $\mathbb{R}^2$ et de ses
deux demi-droites de bord.}

\pagerange{55}{55}
\section*{VII. Les déterminants gradués (page 55)}

Soit $k$ un anneau commutatif, $\mathcal{V}(k)$ le groupoïde des $k$-modules
libres de type fini et de leurs isomorphismes, réunion disjointe des
$\mathcal{V}_n(k)$ des modules de rang $n$, et muni de la somme directe. Soit
$\mathrm{Pic}^{\mathbb{Z}}(k)$ la catégorie de Picard des $k$-modules libres
de rang $1$ gradués par $\mathbb{Z}$, avec pour contrainte de commutativité
l'isomorphisme de Koszul
$\ell \otimes m \mapsto (-1)^{\deg \ell \cdot \deg m}\, m \otimes
\ell$.\footnote{La page note cette catégorie $\mathbb{P}(k)$, et dit
$\mathcal{V}$ « catégorie additive (graduée) » : avec les seuls
isomorphismes pour morphismes, c'est un groupoïde monoïdal symétrique pour
$\oplus$. Le titre, « Théorie des déterminants gradués », est souligné ; un
premier « déterminants » y est barré, avec au-dessus un mot lu sans
certitude « modules ».} Une \emph{théorie des déterminants gradués} est un
foncteur $\Delta : \mathcal{V}(k) \to \mathrm{Pic}^{\mathbb{Z}}(k)$ muni
\begin{itemize}
  \item[a)] d'isomorphismes naturels
  $\varphi_{V,V'} : \Delta(V \oplus V') \simeq \Delta(V) \otimes \Delta(V')$
  qui en font un foncteur monoïdal, d'où $\Delta(0) \simeq \mathbb{1}$ ;
  \item[b)] pour $V$ de rang $1$, d'un isomorphisme $\Delta(V) \simeq V$,
  placé en degré $1$ ;
  \item[c)] de la condition suivante : si $W_2$ et $W'_2$ sont deux
  supplémentaires libres d'un sous-module $W_1 \subset V$, avec
  $\alpha : W_1 \oplus W_2 \simeq V$, $\alpha' : W_1 \oplus W'_2 \simeq V$ et
  $\lambda : W_2 \simeq W'_2$ l'isomorphisme induit par $V/W_1$, le
  diagramme
\end{itemize}

\begin{tikzcd}
\Delta(W_1 \oplus W_2) \arrow[r, "\Delta(\alpha)"] \arrow[d, "{\varphi_{W_1,W_2}}"'] & \Delta(V) & \Delta(W_1 \oplus W'_2) \arrow[l, "\Delta(\alpha')"'] \arrow[d, "{\varphi_{W_1,W'_2}}"] \\
\Delta(W_1) \otimes \Delta(W_2) & & \Delta(W_1) \otimes \Delta(W'_2) \arrow[ll, "{\mathrm{id} \otimes \Delta(\lambda)}"]
\end{tikzcd}

est commutatif.\footnote{La page écrit $\varphi_{W_1, W_1}$ sur la flèche de
gauche, lapsus pour $\varphi_{W_1, W_2}$, et marque les quatre flèches du
signe $\simeq$.} Ce qui équivaut, dit la page, à demander que tout
automorphisme de $V$ qui induit l'identité sur $W_1$ et sur $V/W_1$ agisse
trivialement sur $\Delta(V)$.

La page s'arrête sur les axiomes. Le théorème qu'ils préparent, analogue à
celui de la page 49, n'est pas écrit. L'objet visé est le déterminant gradué
$V \mapsto (\Lambda^{\operatorname{rg} V} V, \operatorname{rg} V)$ ; la
condition c) dit que les transvections agissent trivialement, et sur un
corps, où $\mathrm{GL}_n$ est engendré par les transvections et les matrices
diagonales, l'argument de la page 51 transpose : le caractère par lequel
$\mathrm{GL}_n(k)$ opère sur $\Delta(k^n)$ est forcé d'être le déterminant.
Sur un anneau quelconque, la question touche au groupe $SK_1(k)$, et la page
ne la pose pas.\footnote{Ce paragraphe est le nôtre. Les références sont
Knudsen et Mumford, « The projectivity of the moduli space of stable curves
I : preliminaries on det and Div », Math. Scand. 39 (1976), qui
introduisent le déterminant gradué et ses signes, et Deligne, « Le
déterminant de la cohomologie » (1987), qui en axiomatise la notion ; les
catégories de Picard sont de Deligne (SGA 4, exposé XVIII, 1973), et la thèse
de Hoàng Xuân Sính sur les Gr-catégories (1975) a été dirigée par
Grothendieck. Rien sur la page ne cite ces travaux.}

La règle des signes de la page 53, $\omega' \times \omega =
(-1)^{\dim V \cdot \dim V'}\, \omega \times \omega'$, est la trace de celle-ci :
sur $\mathbb{R}$, l'ensemble des orientations de $V$ est celui des
composantes connexes de $\Delta(V) \smallsetminus \{0\}$, et la contrainte de
Koszul y devient le signe de la page 53. Le rapprochement des deux dernières
sections est le nôtre ; les pages ne le font pas.

\end{document}
